\section{Omitted Details and Proofs for Section \ref{sec:background}}


\subsection{Construction of the Free Algebra for a Signature}\label{sec:app:free-signature-monad}

In this section, we give an explicit construction for the functor that takes a
set $A$ to the free $\Sigma$-algebra $F_\Sigma(A)$ on $A$.
%
We seek to solve the equation
%
\[ F_\Sigma(A) \cong A + \Sigma(F_\Sigma(A)), \]
%
i.e.,
%
\[ F_\Sigma(A) \cong A + \Sigma_{\oper: \Op}, 
    (\tau_\oper \to F_\Sigma(A)) \times G_\oper(\later F_\Sigma(A)). \]
%
First, consider the parameterized inductive type 
%
\[ P(X) := \mu Y. (A + \Sigma_{\oper : \Op}, (\tau_\oper \to Y) \times G_\oper(X)). \]
%
Now, using guarded recursion, define $F_\Sigma(A)$ to be the unique solution to
the guarded-recursive equation
%
\[ F_\Sigma(A) \cong P(\later F_\Sigma(A)). \]
%
Then by the definitions of least fixpoint and guarded fixpoint, we have that
$F_\Sigma(A)$ is a solution to the original equation.

\section{Omitted Details and Proofs for Section \ref{sec:operational-semantics}}

\begin{comment}
\begin{lemma}
    Let $\hat{N} : \Sigmahat(\term{A})$. We have 
    %
    \[ \iota(f (\Sigmahat(E[\cdot])(\hat{N})) ) = (\eta' \circ E[\cdot])^{\ddag} (\iota(f(\hat{N}))). \]
\end{lemma}
\begin{proof}
    We proceed by cases on $\hat{N}$. If $\hat{N} = \cpure{M}$, then we have
    \begin{align*}
        \iota(f (\Sigmahat(E[\cdot])(\cpure{M})) )
        &= \iota(f (\cpure{E[M]}) ) \\
        &= \iota(\eta(E[M])) \\
        &= \eta'(E[M]),
    \end{align*}
    and
    \begin{align*}
        (\eta' \circ E[\cdot])^{\ddag} (\iota(f(\cpure{M})))
        &= (\eta' \circ E[\cdot])^{\ddag} (\iota(\eta(M))) \\
        &= (\eta' \circ E[\cdot])^{\ddag} (\eta'(M)) \\
        &= (\eta' \circ E[\cdot])(M) \\
        &= \eta'(E[M]).
    \end{align*}

    If $\hat{N} = \ceff{(\oper, u)}$, then we have
    \begin{align*}
        \iota(f (\Sigmahat(E[\cdot])(\ceff{(\oper, u)})) )
        &= \iota(f (\ceff{(\oper, E[\cdot] \circ u)}) ) \\
        &= \iota(T_{\mho}(\term{A}).\oper(\eta \circ E[\cdot] \circ u)) \\
        &= T_{\mho sb}(\term{A}).\oper(\iota \circ \eta \circ E[\cdot] \circ u) \\
        &= T_{\mho sb}(\term{A}).\oper(\eta' \circ E[\cdot] \circ u)
    \end{align*}
    and
    \begin{align*}
        (\eta' \circ E[\cdot])^{\ddag} (\iota(f(\ceff{(\oper, u)})))
        &= (\eta' \circ E[\cdot])^{\ddag} (\iota(T_{\mho}(\term{A}).\oper(\eta \circ u))) \\
        &= (\eta' \circ E[\cdot])^{\ddag} (T_{\mho sb}(\term{A}).\oper(\iota \circ \eta \circ u)) \\
        &= (\eta' \circ E[\cdot])^{\ddag} (T_{\mho sb}(\term{A}).\oper(\eta' \circ u)) \\
        &= (T_{\mho sb}(\term{A}).\oper(((\eta' \circ E[\cdot])^{\ddag}) \circ \eta' \circ u)) \\
        &= (T_{\mho sb}(\term{A}).\oper(((\eta' \circ E[\cdot])) \circ u)). \\
    \end{align*}

\end{proof}
\end{comment}


Notation in this section:
%
\begin{itemize}
    \item We let $\eta$ denote $\Tsb.\eta$.
    %
    \item We let $(-)^\dag$ denote Kleisli extension for $\Tsb$.
    %
    \item $f : (\term{A} + \Sigma_\mho(\term{A})) \to \Tsb(\term{A})$ is defined by
    cases, mapping $\cpure{M}$ to $\eta^{\Tsb}(M)$, and mapping $\ceff{(\oper, u, g)}$ 
    to $\Tsb(\term{A}).\oper(\eta \circ u, (\later \eta) \circ g)$.
    (Note that $u : \tau_\oper \to \term{A}$ and $g : \gamma_\oper \to \later (\term A)$.)
    %
    % \item $\iota : T_{\mho}(X) \to T_{\mho sb}(X)$ is the inclusion of the
    % free model monad $T_{\mho}$ into $T_{\mho sb}$. 
\end{itemize}

To begin, we define a function that ``runs'' an evaluation context $E$.
Specifically, if $E : A \multimap B$, then we define $\run{E} : \Tsb(\term A)
\to \Tsb(\val B)$ as
%
\[ \run{E} := (\lambda M. \run{(E[M])})^\dag, \]
%
where $\dag$ is the Kleisli extension operator for the monad $\Tsb$.

Then we have the following lemma, which says that running commutes with filling
an evaluation context:
%
\begin{lemma}[Run is a Homomorphism]\label{lem:run-fill}
  $\run{(E[M])} = (\run{E})(\run{M})$.
\end{lemma}
\begin{proof}
    We want to show that $\run{(E[M])} = (\run{E})(\run{M})$. We proceed by
    guarded recursion and cases on whether $M$ is a value, or $M$ steps.

    If $M$ is a value $V$, then the result is straightforward. 
    % by definition of $\run$ on terms, we have $\run{V} = \eta\, V$.

    Now suppose $M \stepsin{k} \hat{N}$, where $\hat{N} : \Sigmahat(\term A)$.
    Then we have, by definition of $\run$:

    \begin{equation} 
        \run{M} = \delta_{\sty} (\run^{\dag} {(f(\hat{N}))}).
    \end{equation}

    We also have $E[M] \stepsin{k} \Sigmahat(E[\cdot])(\hat{N})$. Thus, by the
    definition of $\run$, we have:
    \begin{equation} 
        \run{(E[M])} = \delta_{\sty} (\run^{\dag} {(f (\Sigmahat(E[\cdot])(\hat{N})) )}).
    \end{equation}

    We continue by cases on $\hat{N}$. If $\hat{N} = \cpure{M}$, then we have
    \begin{align*}
        \run{E[M]}
        &= \delta_{\sty} [\run^{\dag} {(f (\Sigmahat(E[\cdot])(\cpure{M})) )}] \\
        &= \delta_{\sty} [\run^{\dag} {(f (\cpure(E[M])))}] \\
        &= \delta_{\sty} [\run^{\dag} {\eta (E[M])}] \\
        &= \delta_{\sty} [\run E[M]]
    \end{align*}

    On the other hand, we have
    \begin{align*}
        (\run E)(\run M) 
        &= (\lambda V. \run{(E[V])})^{\dag} (\delta_{\sty} (\run^{\dag} {(f(\cpure{M}))}) ) \\
        &= \delta_{\sty} [(\lambda V. \run{(E[V])})^{\dag} (\run^{\dag} {(f(\cpure{M}))})] \\
        &= \delta_{\sty} [(\lambda V. \run{(E[V])})^{\dag} (\run^{\dag} {\eta M})] \\
        &= \delta_{\sty} [(\lambda V. \run{(E[V])})^{\dag} (\run M)] \\
        &= \delta_{\sty} [(\run E) (\run M)]
    \end{align*}
    %
    Since both sides are guarded by a $\delta_{\sty}$, it suffices to show that
    the insides are equal later. This follows from our \lob~induction
    hypothesis.

    Now suppose $\hat{N} = \ceff{(\oper, u)}$. Then for the LHS, we have
    %
    \begin{align*}
        \run{(E[M])} 
        &= \delta_{\sty} [\run^{\dag} {(f (\Sigmahat(E[\cdot])(\ceff{(\oper, u)})) )}] \\
        &= \delta_{\sty} [\run^{\dag} {(f (\ceff{(\oper, E[\cdot] \circ u)}) )}] \\
        &= \delta_{\sty} [\run^{\dag} {(\Tsb(\term{A}).\oper(\eta \circ E[\cdot] \circ u) )}] \\
        &= \delta_{\sty} [\Tsb(\term{A}).\oper(\run^{\dag} \circ \eta \circ E[\cdot] \circ u)] \\
        &= \delta_{\sty} [\Tsb(\term{A}).\oper(\run \circ E[\cdot] \circ u)] \\
        &= \delta_{\sty} [\Tsb(\term{A}).\oper(\lambda w. \run (E[u(w)]))] \\
    \end{align*}

    For the RHS, we have
    %
    \begin{align*}
        (\run E)(\run M) 
        &= (\lambda V. \run{(E[V])})^{\dag} (\delta_{\sty} (\run^{\dag} {(f(\ceff{(\oper, u)}))}) ) \\
        &= \delta_{\sty} [(\lambda V. \run{(E[V])})^{\dag} (\run^{\dag} {(f(\ceff{(\oper, u)}))})] \\
        &= \delta_{\sty} [(\lambda V. \run{(E[V])})^{\dag} (\run^{\dag} { \Tsb(\term{A}).\oper(\eta \circ u) })] \\
        &= \delta_{\sty} [(\lambda V. \run{(E[V])})^{\dag} (\Tsb(\term{A}).\oper(\run^{\dag} \circ \eta \circ u))] \\
        &= \delta_{\sty} [(\lambda V. \run{(E[V])})^{\dag} (\Tsb(\term{A}).\oper(\run \circ u))] \\
        &= \delta_{\sty} [\Tsb(\term{A}).\oper((\lambda V. \run{(E[V])})^{\dag} \circ \run \circ u)] \\
        &= \delta_{\sty} [\Tsb(\term{A}).\oper(\lambda w. (\run E) (\run (u(w))))] \\
    \end{align*}
    %
    As before, since both sides are guarded by a $\delta_{\sty}$, it suffices to
    show that the insides are equal later, which follows from our \lob~induction
    hypothesis.

    % Thus, it will suffice to show that
    % \[ 
    %   \later[ (\run \circ E[\cdot])^{\dag} (\iota(f(\hat{N}))) 
    %     = ((\lambda V. \run{(E[V])})^{\dag} \circ \run)^{\dag} {(\iota(f(\hat{N})))} ]. 
    % \]


\end{proof}

\section{Omitted Details and Proofs for Section \ref{sec:denotational-semantics}}

\subsection{Action of the Signature Endofunctor on Relations}

We describe the action of the endofunctor for the algerbaic signature $\Sigma$
on predomain relations.

Let $c : A \rel A'$ be a predomain relation. The action of the endofunctor
$\Sigma$ on $c$, written $\Sigma(c)$, is a predomain relation between
$\Sigma(A)$ and $\Sigma(A')$, defined as follows:
%
First, note that any predomain relation $c : A \rel A'$ extends to a predomain
relation $\tau \to c$ between $\tau \to A$ and $\tau \to A'$ defined pointwise,
i.e., $f \binrel{(\tau \to c)} f'$ iff for all $s : \tau$ we have $(f s)
\binrel{c} (f' s)$.
%
Similarly, given $c : A \rel A'$ we define $\later c$ such that $\tilde{x}
\binrel{(\later c)} \tilde{y}$ iff after one step passes, the left and right
sides (now in $A$ and $A'$ respectively) are related in $c$.
%
We now define $(\oper, (u, g)) \binrel{(\Sigma(c))} (\oper', (u', g'))$ iff
$\oper = \oper'$ and $u \binrel{(\tau_\oper \to c)} u'$ and 
$g \binrel{(\gamma_\oper \to \later c)} g'$.

\begin{comment}
\subsection{Generalizing Error Domains}

In this section, we discuss the guarded denotational model for gradual typing.

We begin by reviewing the key ingredients in the denotational semantics of
Giovannini et al. \cite{gtt-sgdt}. For further discussion, see their paper. For
the sake of brevity, we will refer to their model as the \emph{prior guarded
model} of gradual typing.

The prior model is based on call-by-push-value \cite{levy99}, a categorical
model that refines Moggi's monadic semantics \cite{MOGGI199155}. In
call-by-push-value, there is a category of value types $\calV$ and pure
morphisms $f : A \to A'$, and a category $\calE$ of computation types $\phi : B
\multimap B'$. A pair of adjoint functors $F \vdash U$ mediates between the two
categories, where $F : \calV \to \calE$ takes $A$ to the type $F A$ of
computations that return values in $A$, and $U : \calE \to \calV$ takes $B$ to
the type $U B$ of suspended computations of type $B$.

In the prior model, the value types are modeled by sets equipped with a partial
ordering relation that models type precision on closed values. However, since
step-insensitive reasoning does not play well with transitivity, each type also
has an additional relation called \textbf{weak bisimilarity}. They call such a
set a \textbf{predomain}. Morphisms of predomains are functions on the
underlying sets that respect both of the relations.
%
\begin{ifextended}
The precise definitions are given below:
%
\begin{definition}
A \textbf{predomain} $A$ consists of a set $A$ along with two relations:
\begin{itemize}
    \item A partial order $\ltdyn_A$.
    \item A reflexive, symmetric ``bisimilarity'' relation $\bisim_A$.
\end{itemize}

A \textbf{morphism of predomains} $f : A \to A'$ is a function between the
underlying sets that is both \emph{monotone} (if $x \ltdyn_A x'$, then $f(x)
\ltdyn_{A'} f(x')$), and \emph{bisimilarity-preserving} (if $x \bisim_A x'$,
then $f(x) \bisim_{A'} f(x')$).
\end{definition}
\end{ifextended}
%
Computation types are equipped with the same two relations, in addition to a
special distinguished error element that is the least element in the ordering,
as well as a distinguished morphism of predomains $\theta_B : \later B \to B$.
%
The precise definitions are given below:
%
\begin{definition}
An \textbf{error domain} $B$ consists of a predomain $B$ along with the
following data:
\begin{itemize}
    \item A distinguished ``error'' element $\mho_B \in B$ such that $\mho
    \ltdyn_B x$ for all $x$
    \item A morphism of predomains $\theta_B \colon \later B \to B$
    \item For all $x, y : B$ with $x \bisim_B y$, we have $\theta_B(\nxt\, x)
    \bisim_B y$
\end{itemize}
\end{definition}

A \textbf{morphism of error domains} is a morphism of the underlying predomains
that preserves the error element and the $\theta$ map.
%
For an error domain $B$, we define the predomain morphism $\delta_B := \theta_B
\circ \nxt$.

In the prior model, type precision is modeled as a kind of heterogeneous
relation on the denoted predomains. More concretely, A \textbf{predomain
relation} $c : A \rel A'$ is a relation between $A$ and $A'$ that is upward
closed in the ordering on $A'$ and downward closed in the ordering on $A$
(no condition is required involving the bisimilarity relations).
%
\begin{definition}
  A \textbf{predomain relation} $c : A \rel A'$ is a relation $c$ between the
  underlying sets of $A$ and $A'$ satisfying:
  \begin{enumerate}
    \item (upward closure): For all $x \in A$ and $y, y' \in A'$, if $x
    \binrel{c} y$ and $y \ltdyn_{A'} y'$, then $x \binrel{c} y'$.
    \item (downward closure): For all $x, x' \in A$ and $y \in A'$, if $x'
    \ltdyn_{A} x$ and $x \binrel{c} y$, then $x' \binrel{c} y$.
  \end{enumerate}
\end{definition}
%
The ``reflexive'' predomain relation is written $r(A)$ and is given by the
ordering $\ltdyn_A$. The relational composition of two predomain relations is
also a predomain relation, and $r(A)$ is the identity for this composition.

Error domain relations are predomain relations that are additionally a
congruence with respect to the operations $\theta$ and $\mho$.
%
\begin{definition}
  An \textbf{error domain relation} $d : B \rel B'$ is a relation $d$ on the
  underlying predomains satisfying
  \begin{enumerate}
     \item ($\mho$ congruence): $\mho_B \binrel{d} \mho_{B'}$.
     \item ($\theta$ congruence): For all $\tilde{x}$ in $\laterhs B$ and
     $\tilde{y} \in \laterhs B'$, if $\laterhs_t (\tilde{x}_t \binrel{d}
     \tilde{y}_t)$, then $(\theta_B (\tilde{x}) \binrel{d} \theta_{B'}
     (\tilde{y}))$. 
  \end{enumerate}
\end{definition}
Note that because $\mho_B$ is a least element this implies the
stronger property $\mho_B \binrel{d} y$ for any term $y: B'$.
%
There is again a reflexive error domain relation $r(B)$ given by the ordering.
The relational composition of two error domain relations in general is not an
error domain relation ($\theta$ congruence in particular is not preserved), but
we instead use an alternative ``free composition'' $dd' : B \rel B''$ that is
inductively generated from the relational composition.

Lastly, term precision is modeled via the notion of \emph{square}, which
captures the logical-relations style relationship between the functions denoted
by two terms.
%
More concretely:
%
\begin{definition}
Let $f : A_i \to A_o$ and $g : A_i' \to A_o'$ be predomain morphisms, and let
$c_i : A_i \rel A_i'$ and $c_o : A_o \rel A_o'$ be predomain relations. There is
a \textbf{predomain square} $f \ltsq{c_i}{c_o} g$ if for all $x \binrel{c_i} y$,
we have $f(x) \binrel{c_o} g(y)$. 

Likewise, an \textbf{error domain square} $\phi \ltsq{d_i}{d_o} \psi$ is defined
similarly, with no extra interaction with the algebraic structure.
\end{definition}
%
\end{comment}

\begin{comment}
\subsection{Perturbations}

\begin{definition}[Value and Computation Objects and Morphisms]
    A \textbf{value object} consists of a predomain $A$, a monoid $M_A$ and a
    homomorphism of monoids $i_A \colon M_A \to \morbisimid{A} := \{ f : A \to A
    \mid f \bisim \id_A \}$.

    A \textbf{computation object} consists of an error domain $B$, a monoid
    $M_B$ and a homomorphism of monoids $i_B \colon M_B \to \morbisimid{B}$.

    We define a \textbf{value morphism} to be a morphism of the underlying
    predomains, and a \textbf{computation morphism} to be a morphism of the
    underlying error domains.
\end{definition}

\begin{definition}[Push-Pull Structure]
    Let $(A, M_A, i_A)$ and $(A', M_{A'}, i_{A'})$ be value objects, and let $c
    : A \rel A'$ be a predomain relation. A \textbf{push-pull structure} for $c$
    (with respect to $M_A$ and $M_{A'}$) consists of monoid homomorphisms $\push
    : M_A \to M_A'$ and $\pull : M_A' \to M_A$ such that for all $m_A \in M_A$
    there is a square $i_A(m_A) \ltsq{c}{c} i_{A'}(\push\, m_A)$ and for all
    $m_{A'} \in M_{A'}$ there is a square $i_A(\pull\, m_{A'}) \ltsq{c}{c}
    i_{A'}(m_{A'})$. A push-pull structure for error domain relations is defined
    similarly.
\end{definition}

\begin{definition}[Value and Computation Relations]
    A \textbf{value relation} between value objects $(A, M_A, i_A)$ and $(A',
    M_{A'}, i_{A'})$ consists of a predomain relation $c : A \rel A'$ that has a
    push-pull structure, is quasi-left-representable by a morphism $e_c : A \to
    A'$, and is such that $\li c$ is quasi-right-representable by an error
    domain morphism $p_c : \li A' \to \li A$.

    A \textbf{computation relation} between computation objects $(B, M_B, i_B)$
    and $(B', M_{B'}, i_{B'})$ consists of an error domain relation $d$ that has
    a push-pull structure, is quasi-right-representable by a morphism $p_d : B'
    \to B$, and is such that $Ud$ is quasi-left-representable by a morphism $e_d
    : UB \to UB'$
\end{definition}
\end{comment}

\subsection{Adjusting the Term Semantics for Casts}
We provide more details on the change we make to the denotational semantics so
that the denotational and operational semantics are in lock-step up to the
presence of bureaucratic delays in the operational semantics. For an overview,
see Section \ref{sec:perturbations} in the paper.

We begin by recalling that term precision $M \ltdyn N : A$ is modeled in the
semantics of Giovannini et al. by $\sem{M} \ltbisim \sem{N}$ where
%
\[ f \ltbisim g := \exists f', g'. f \bisim f' \ltdyn g' \bisim g. \]
%
That is, we do not require that $f$ and $g$ are in lock-step; we are allowed to
use any $f'$ bisimilar to $f$ and $g'$ bisimilar to $g$ such that $f'$ is
lock-step related to $g'$ in the error ordering.
%
This is useful because the downcast associated to the composition of type
precision $c : A \ltdyn A'$ and $c' : A' \ltdyn A''$ is a morphism from
$F(\sem{A''})$ to $F(\sem{A})$, but in the prior model it is not simply the
composition of the downcast morphisms for $c$ and $c'$. Instead, an additional
perturbation is inserted in order to make the relational squares line up in the
proof of quasi-representability. More details can be found in the Appendix of
their paper \cite{gtt-sgdt-extended-version}.

On the other hand, observe that we could instead define the downcast morphism
for the composition $cc'$ to be the composition of the corresponding downcasts,
without the additional perturbation, and then note that the resulting morphism
is \emph{weakly bisimilar} to the version with the additional perturbation. This
is sufficient for graduality, since the additional perturbation will be absorbed
by the weak bisimilarity in the definition of $f \ltbisim g$.

% Using the above definition for term precision, we see that the additional
% perturbations in the cast corresponding to a composition will be absorbed by the
% weak bisimilarity.

With this observation in mind, we augment the denotation of type precision in
the prior denotational model by associating to each value relation $c$ in the
semantics two morphisms: a ``relational'' embedding morphism $e$ such that $c$
is quasi-left-representable by $e$, and a ``term'' embedding morphism $e' \bisim
e$. We similarly associate two projection morphisms to each computation relation
$d$. We use the term versions of the embedding and projection morphisms as our
denotations of upcasts and downcasts, and we use the relational versions
internally in the proof of graduality.
%
The result is that our term semantics for casts no longer involves unnecessary
perturbations that were inserted to make the graduality proof work.



\section{Omitted Details and Proofs for Section \ref{sec:relating-the-semantics}}

\subsection{Relational Lifting that Ignores Bureaucratic Steps}

  \begin{mathpar}
    \inferrule*[right=RetRet]
    { x \binrel{R} y }
    { \eta^{T_{\mho s}}(x) \binrel{(\bbt(R))} \eta^{T_{\mho sb}}(y) }

    \inferrule*[right=OpOp]
    { \oper \in \calA_{\mho s}.\Op 
      \and \forall s.\, (u\, s) \binrel{(\bbt{R})} (u'\, s) \\
      \and \forall s.\, (g\, s) \binrel{(\later(\bbt{R}))} (g'\, s)  
    }
    { (T_{\mho s}(X).\oper(u, g)) \binrel{(\bbt{R})} (T_{\mho sb}(Y).\oper(u', g')) }

    \inferrule*[right=IdStep]
    { \bar{x} \binrel{(\bbt(R))} \bar{y} }
    { \bar{x} \binrel{(\bbt(R))} \delta_b(\bar{y}) }

  \end{mathpar}

\subsection{Properties of the Logical Relation}

First, we prove two anti-reduction lemmas, one involving bureaucratic steps and
the other involving true steps.

%Proof of the anti-reduction lemma (Lemma \ref{lem:anti-reduction}):
\begin{lemma}[Anti-Reduction]\label{lem:anti-reduction}

  Let $\cdot \vdash M : A$, and suppose $M \stepsin{0} \cpure{M'}$. If
  $\erel{A}{\xbar}{M'}$, then $\erel{A}{\xbar}{M}$.

\end{lemma}
\begin{proof}
    Let $\cdot \vdash M : A$, where $M \stepsin{0} \cpure{M'}$, and suppose
    $\erel{A}{\xbar}{M'}$. We will show that $\erel{A}{\xbar}{M}$.

    % Let $\sty$ be $\bstep$ or $\tstep$ according to whether $k$ is $0$ or $1$,
    % respectively.

    By the definition of $\run$, we have that
    \begin{align*}
        \run{M} &= \delta_{\bstep} (\run^{\dag} {(f(\cpure{M'}))}) \\
                &= \delta_{\bstep} (\run^{\dag} {\eta(M')}) \\
                &= \delta_{\bstep} (\run{M'}).
    \end{align*}

    Since $\erel{A}{\xbar}{M'}$, we have that $\liftvrel{A}{\xbar}{(\run M')}$.
    By definition of the relational lifting, we have that
    $\liftvrel{A}{\xbar}{\delta_{\bstep}(\run{M'})}$. Thus, we have
    $\liftvrel{A}{\xbar}{\run{M}}$, which means that $\erel{A}{\xbar}{M}$.

\end{proof}

\begin{lemma}[Anti-Reduction for True Steps]\label{lem:anti-reduction-true}
    Let $\cdot \vdash M : A$, and suppose $M \stepsin{1} \cpure{M'}$. Then for
    all $\tilde{x} : \later \Ts(\sem{A})$, if $\later_t
    (\erel{A}{\tilde{x}_t}{M'})$, then $\erel{A}{\theta(\tilde{x})}{M}$.
\end{lemma}
\begin{proof}
    Suppose $\later_t (\erel{A}{\tilde{x}_t}{M'})$. We will show that
    $\erel{A}{\theta(\tilde{x})}{M}$. By the definition of $\run$, we have that
    \begin{align*}
        \run{M} &= \delta_{\tstep} (\run^{\dag} {(f(\cpure{M'}))}) \\
                &= \delta_{\tstep} (\run^{\dag} {\eta(M')}) \\
                &= \delta_{\tstep} (\run{M'}).
    \end{align*}

    It suffices to show $\liftvrel{A}{\theta(\tilde{x})}{\run{M}}$,
    i.e., $\liftvrel{A}{\theta(\tilde{x})}{\delta_{\tstep}(\run{M'})}$.
    Now using the definition of the relational lifting, it suffices to show that
    \[ \later_t(\liftvrel{A}{\tilde{x}_t}{(\nxt (\run{M'}))_t}), \]
    i.e.,
    \[ \later_t(\liftvrel{A}{\tilde{x}_t}{(\run{M'})}). \]
    But this is by definition the same as
    \[ \later_t(\erel{A}{\tilde{x}_t}{M'}), \]
    which is our assumption.
\end{proof}


Next, we have a semantic bind lemma for the logical relation, which allows us to
replace arbitrary terms in an evaluation context with values.

%Proof of the Semantic Bind Lemma:
\begin{lemma}[Monadic Bind]\label{lem:monadic-bind}
    % Suppose for all values $V, W$ that $\vrel{A}{V}{W}$ implies
  % $\erel{B}{\phi(V)}{(E[W])}$. Then, for all $\erel{A}{\bar{x}}{N}$, we have
  % $\erel{B}{(\phi(\bar{x}))}{(E[N])}$.
  Let $\phi : T_{\mho s}(\sem{A}) \to T_{\mho s}(\sem{B})$, and let $\bullet : A \vdash E : B$.
  %
  Suppose for all values $x, V$ that $\vrel{A}{x}{V}$ implies $\erel{B}{\phi(T_{\mho s}.\eta(x))}{(E[V])}$.
  %
  Then for all $\bar{x} \in T_{\mho s}(\sem{A})$ and terms $N$ of type $A$ such that
  $\erel{A}{\bar{x}}{N}$, we have $\erel{B}{(\phi(\bar{x}))}{(E[N])}$
\end{lemma}
\begin{proof}
    Follows from Lemma \ref{lem:run-fill} and the universal property of the
    relational lifting $\bbt(R)$.
\end{proof}

\begin{comment}
\begin{proof}
    Suppose for all values $x, V$ that $\vrel{A}{x}{V}$ implies
    $\erel{B}{\phi(T_{\mho s}.\eta(x))}{(E[V])}$. Further, suppose that
    $\erel{A}{\bar{x}}{N}$.
    %
    We will show that $\erel{B}{(\phi(\bar{x}))}{(E[N])}$, which by definition
    is $\liftvrel{B}{\phi(\bar{x})}{(\run (E[N]))}$

    The proof is by guarded recursion and induction over the proof that
    $\erel{A}{\bar{x}}{N}$, since this is defined as
    $\liftvrel{A}{\bar{x}}{(\run N)}$. We have the following cases:

    \begin{itemize}
        \item Case \textsc{RetRet}: This means that $\bar{x} = \eta\, x$ for
        some $x \in \sem{A}$ and $\run N = \eta'\, V$ for some $V \in \val{A}$
        where $\vrel{A}{x}{V}$. (Here, $\eta$ is the unit for $T_{\mho s}$ and
        $\eta'$ is the unit for $T_{\mho sb}$.) Then by our assumption, we have
        $\erel{B}{\phi(x)}{(E[V])}$, as required.

        \item Case \textsc{OpOp}: In this case, we have
        %
        $\bar{x} = T_{\mho s}(\sem{A}).\oper(u, g)$ for some $u$ and $g$.
        %
        We also know that
        %
        $\run N = (T_{\mho sb}(\val A).\oper(u', g'))$
        %
        for some $u'$ and $g'$, where for all $s$, we have $\liftvrel{A}{(u\,
        s)}{(u'\, s)}$ and $(g\, s) \binrel{(\later(\bbt{\vrel{A}}))} (g'\, s)$.

        Now since $\phi$ is a homomorphism, we have
        %
        \[ \phi(T_{\mho s}(\sem{A}).\oper(u, g)) = 
           T_{\mho s}(\sem{B}).\oper(\phi \circ u, (\later \phi) \circ g).
        \]
        %
        Furthermore, by Lemma \ref{lem:run-fill}, we have $\run (E[N]) = (\run
        E)(\run N)$ where we recall that $\run E$ is defined as $(\lambda V.
        \run{(E[V])})^\dag$. In particular, this means that
        \begin{align*}
            \run (E[N]) &= (\lambda V.\run{(E[V])})^\dag (\run N)  \\
                        &= (\lambda V.\run{(E[V])})^\dag (T_{\mho sb}(\val A).\oper(u', g')) \\
                        &= T_{\mho sb}(\val B).\oper((\run E) \circ u', (\later (\run E)) \circ g')
        \end{align*}

        Thus, using rule $\textsc{OpOp}$ and the induction hypotheses, this case is complete.
        % TODO: more detail?

        
        \item Case \textsc{IdStep}: In this case, we have
        %
        $\run N = \delta_b(\bar{V})$
        %
        for some $\bar{V} \in T_{\mho sb}(\val{A})$, and
        $\liftrel{A}{\bar{x}}{\bar{V}}$.
                
        Then as above, we have
        %
        \begin{align*}
            \run (E[N]) &= (\run E)(\run N) \\
                &= T_{\mho sb}(\val B).\delta_b[(\run E)(\bar{V})] \\
                &= 
        \end{align*} 
        %
        so by rule \textsc{IdStep}, it suffices to show
        %
        \[ \liftrel{B}{\phi(\bar{x})}{(\run E)(\bar{V})} \]
    \end{itemize}
\end{proof}
\end{comment}

\subsection{Proof of the Fundamental Lemma for the Logical Relation}

Towards proving the fundamental lemma, we prove a series of compatibility
lemmas, one for each term constructor in the cast calculus.

\begin{lemma}[Logical Relation for Natural Number Values]\label{lem:lr-nat}
    Let $n \in \mathbb{N}$, and let $V \in \val{\nat}$.
    Then if $\vrel{\nat}{n}{V}$, then $\vrel{\nat}{(\suc\, n)}{(\suc\, V)}$
\end{lemma}
\begin{proof}
    If $\vrel{\nat}{n}{V}$, then by definition we have $V = \suc^n(\zro)$. This
    means that $\suc\,V = \suc^{n+1}(\zro)$, i.e., $\vrel{\nat}{(\suc\,
    n)}{(\suc\, V)}$.
\end{proof}

\begin{lemma}[Compatibility for Zero]
    We have $\simrel{\Gamma}{\nat}{\sem{\zro}}{\zro}$.
\end{lemma}
\begin{proof}
    Follows by definition of $\vrel{\nat}{}{}$.
\end{proof}


\begin{lemma}[Compatibility for Successor]
    If $\simrel{\Gamma}{\nat}{\sem{M}}{M}$, then
    $\simrel{\Gamma}{\nat}{\sem{\suc\, M}}{\suc\, M}$.
\end{lemma}
\begin{proof}
    Let $(\grel{\Gamma} {\gambar} {\gamma'})$. We need to show 
    %
    \[ \erel{\nat} {(\sem{\suc\, M}\, \gambar)} {\suc\, M[\gamma']}. \]
    %
    We have $\sem{\suc\, M}\, \gambar = (\lambda x.\eta(\suc\, x))^\dag(\sem{M}\, \gambar)$
    
    By the bind lemma, it suffices to show that (1) $\erel{\nat}{\sem{M}\,
    \gambar}{M[\gamma']}$ and (2) for all $x, V$ such that $\vrel{\nat}{x}{V}$,
    we have $\erel{\nat}{\eta(\suc\,x)}{\suc\,V}$.
    Note that (1) follows from our assumption.

    To show (2), we note that both are values, so it suffices to show that
    $\vrel{\nat}{\suc\, x}{\suc\, V}$. Then by Lemma \ref{lem:lr-nat}, it
    suffices to show $\vrel{\nat}{x}{V}$, which is what we have assumed.    
\end{proof}

\begin{lemma}[Compatibility for Variables]
    For all $\Gamma_1, \Gamma_2$, we have $\simrel{\Gamma_1, x : A, \Gamma_2}{A}{\sem{x}}{x}$.
\end{lemma}
\begin{proof}
    Let $\grel{\Gamma_1, x : A, \Gamma_2}{\gambar}{\gamma'}$. We need to show
    \[ \erel{A}{\sem{x}\, \gambar}{x[\gamma']}.  \]
    We have $\sem{x}\, \gambar = \eta(\gambar(x))$, and $x[\gamma'] = \gamma'(x)$.
    Since both are values, it suffices to show
    \[ \vrel{A}{\gambar(x)}{\gamma'(x)}, \]
    which follows from the assumption that $\gambar$ is related to $\gamma'$.
\end{proof}

\begin{lemma}[Compatibility for Lambda]
    Let $\fbar : \sem{\Gamma} \times \sem{A_i} \to T\sem{A_o}$.
    %
    Suppose $\simrel{\Gamma, x : A_i}{A_o}{\fbar}{M}$.
    %
    Then we have $\simrel{\Gamma}{A_i \ra A_o}{\bar{m}}{\lambda x.M}$, where
    $\bar{m}(\gambar) = \eta(\lambda x. \fbar(\gambar, x))$.
\end{lemma}
\begin{proof}
    % \sem{M} : \sem{\Gamma, x : A_i} \to T(\sem{A_o})
    % = \sem{\Gamma} \times \sem{A_i} \to T(\sem{A_o})
    %
    % \sem{\lambda x.M} : \sem{\Gamma} \to T(\sem{A_i \ra A_o})
    % \sem{\lambda x.M} \gambar = \eta( \lambda x. \sem{M}(\gambar , x) )

    Let $\grel{\Gamma}{\gambar}{\gamma'}$. We need to show
    %
    \[ \erel{A_i \ra A_o} {(\bar{m}\, \gambar)} {(\lambda x.M)[\gamma']}. \]
    %
    Since both are values, it suffices to show that
    %
    \[ \vrel {A_i \ra A_o} {(\lambda x. \fbar(\gambar, x))} {(\lambda x.M[\gamma'])}. \]
    %
    Now, by definition of $\vrel{A_i \ra A_o}{}{}$, let $x : \sem{A_i}$ and $V : \val{A_i}$,
    and suppose $\vrel{A_i}{x}{V}$. We will show that
    \[ \erel{A_o} {((\lambda x. \fbar(\gambar, x))\, x)} {((\lambda x.M[\gamma'])\, V)}, \]
    i.e.,
    \[ \erel{A_o} {(\fbar(\gambar, x))} {((\lambda x.M[\gamma'])\, V)}. \]
    Since $((\lambda x.M[\gamma'])\, V) \stepsin{0} (M[\gamma'][V/x])$,
    by anti-reduction it suffices to show
    \[ \erel{A_o} {(\fbar(\gambar, x))} {(M[\gamma'][V/x])}. \]

    Let $\bar{\gamma_1} := (\gambar, x) \in \sem{\Gamma, x : A_i}$.
    Let $\gamma_1' := (\gamma', x \mapsto V)$.
    Then we have $\grel{\Gamma, x : A_i}{\bar{\gamma_1}}{\gamma'}$,
    so by our assumption that $\simrel{\Gamma, x : A_i}{A_o}{\fbar}{M}$,
    we have $\erel{A_o} {\fbar\, \bar{\gamma_1}} {M[\gamma_1']}$.
    This is what we needed to show.

\end{proof}

\begin{lemma}[Compatibility for Function Application]
    Let $\fbar : \sem{\Gamma} \to T\sem{A_i \ra A_o}$ and 
    $\xbar : \sem{\Gamma} \to T\sem{A_i}$.
    %
    Suppose $\simrel{\Gamma}{A_i \ra A_o}{\fbar}{M}$, and
    $\simrel{\Gamma}{A_i}{\xbar}{N}$. Then we have
    $\simrel{\Gamma}{A_o}{\bar{m}}{M\, N}$ where

    $\bar{m}\, \gambar := (\lambda g. g^\dag (\xbar\, \gambar))^\dag (\fbar\, \gambar)$
\end{lemma}
\begin{proof}
    Let $\grel{\Gamma}{\gambar}{\gamma'}$. We need to show
    %
    \[ \erel{A_o} {\bar{m}\, \gambar} {(M[\gamma']\, N[\gamma'])}. \]

    By the Semantic Bind Lemma with
    %
    $\phi(\bar{x_f}) = (\lambda g. g^\dag (\xbar\, \gambar))^\dag (\bar{x_f})
    %
    : \Ts(\sem{A_i \ra A_o}) \to \Ts(A_o)$, and 
    %
    $E = \bullet\, N[\gamma']$, it suffices to show:
    %
    (1) $\erel{A_i \ra A_o}{(\fbar\, \gambar)}{M[\gamma']}$ and
    %
    (2) for all $x_f : \sem{A_i \ra A_o}$ and $\Gamma \vdash V_f : A_i \ra A_o$,
    if $\vrel{A_i \ra A_o}{x_f}{V_f}$, then $\erel{A_o} {(\phi(\eta\, x_f))} {V_f\, (N[\gamma'])}$.

    (1) follows from our assumption that $\simrel{\Gamma}{A_i \ra
    A_o}{\fbar}{M}$.
    %
    To show (2), we again apply the Semantic Bind Lemma. This time, it suffices
    to show that for all $x : \sem{A_i}$ and $\Gamma \vdash V : A_i$, if
    $\vrel{A_i}{x}{V}$, then $\erel{A_o}{x_f\, x}{V_f\, V}$. This follows from
    our assumption that $\vrel{A_i \ra A_o}{x_f}{V_f}$.

    % Note: \phi(\eta\, x_f) = (\lambda f. f^\dag (\sem{N} \gambar))^\dag (\eta\, x_f)
    %                        = x_f^\dag (\sem{N} \gambar) 


\end{proof}

\begin{lemma}[Compatibility for Product Split]
Suppose $\simrel{\Gamma}{A_1 \times A_2}{\sem{M}}{M}$ and 
$\simrel{(\Gamma, x : A_1, y : A_2)}{A}{\sem{N}}{N}$. Then we have
$\simrel {\Gamma} {A} {\sem {M_{\text{split}}}} {{M_{\text{split}}}}$,
where ${M_{\text{split}}} = \textrm{split } (x,y) = M \textrm{ in } N$.
\end{lemma}
\begin{proof}
    Follows by the semantic bind lemma and the given hypotheses.
\end{proof}

\begin{lemma}[Compatibility for Pairing]
Suppose $\simrel{\Gamma}{A_1}{\sem{M_1}}{M_1}$ and $\simrel{\Gamma}{A_2}{\sem{M_2}}{M_2}$.
Then we have $\simrel{\Gamma}{A_1 \times A_2}{\sem{(M_1, M_2)}}{(M_1, M_2)}$.
\end{lemma}
\begin{proof}
    Follows by the semantic bind lemma and the given hypotheses.
\end{proof}


\begin{lemma}[Compatibility for Effects]
    Let $\oper$ be an effect with arity $n$. Suppose for each $i \in \{1, \dots,
    n\}$ we have $\simrel{\Gamma}{A}{\sem{M_i}}{M_i}$. Then we have
    $\simrel{\Gamma}{A}{\sem{\synoper{M}{n}}}{(\synoper{M}{n})}$.
\end{lemma}
\begin{proof}
    Let $(\grel{\Gamma}{\gambar}{\gamma'})$. We need to show that
    %
    \[ \erel{A} {(\sem{\synoper{M}{n}}\, \gambar)} {(\synoper{M}{n})[\gamma']}. \]
    %
    First note that
    \[ \sem{\synoper{M}{n}}\, \gambar = T_{\mho s}.\oper(\lambda i.\, (\sem{M_i}\, \gambar)). \]
    %
    Furthermore, note that $\synopernoidx{M_i[\gamma']}{n}$ steps to
    $\hat{N} := \ceff{(\oper, (\lambda\, (i : [n]).\, M_i[\gamma']))}$, so by definition of
    $\run$, we have
    %
    \begin{align*}
        \run{\synopernoidx{M_i[\gamma']}{n}}
        &= \delta_b (\run^{\dag} {(f(\hat{N}))}) \\
        &= \delta_b (\run^{\dag} \Tsb(\term{A}).\oper(\eta \circ (\lambda\, i.\, M_i[\gamma']))) \\
        &= \delta_b (\Tsb(\term{A}).\oper(\run^{\dag} \circ \eta \circ (\lambda\, i.\, M_i[\gamma']))) \\
        &= \delta_b (\Tsb(\term{A}).\oper(\run \circ (\lambda\, i.\, M_i[\gamma']))) \\
        &= \delta_b (\Tsb.\oper(\lambda i.\, \run(M_i[\gamma']))). 
    \end{align*}

    Now, using rules \textsc{IdStep} and \textsc{OpOp} of the relational
    lifting, it suffices to show that for each $i \in [n]$, we have
    %
    \[ \liftvrel{A}{(\sem{M_i}\, \gambar)}{(\run (M_i[\gamma']))}, \]
    %
    which follows from our assumption.


\end{proof}

\begin{lemma}[Compatibility for Casts]
    Let $c : A \ltdyn A'$, and let $\fbar : \sem{\Gamma} \to T\sem{A}$.
   
    If $\simrel{\Gamma}{A}{\fbar}{M}$, then we have
    $\simrel{\Gamma}{A'}{(T(\sem{\upc c}) \circ \fbar)}{\upc{c} M}$.

    Likewise, let $\gbar : \sem{\Gamma} \to T\sem{A'}$.
    If $\simrel{\Gamma}{A'}{\gbar}{M}$, then we have
    $\simrel{\Gamma}{A}{\sem{\dnc c} \circ \gbar}{\dnc{c} M}$.
\end{lemma}
\begin{proof}
    The proof is by guarded recursion and induction on the type precision
    derivation $c$. Let $(\grel{\Gamma}{\gambar}{\gamma'})$.

    We consider the upcasts first. We need to show
    %
    \[ \erel {A'} {T\sem{\upc c} (\fbar\, \gambar)} {(\upc c M)[\gamma']}. \]
    %
    Since casts are evaluation contexts, it suffices by the semantic bind lemma
    to assume that the inputs to the casts are values.
    Specifically, let $x : \sem{A}$ and $V : \val{A}$, and suppose that $\vrel{A}{x}{V}$.
    We will show that

    \[ \erel{A'}{(\eta\, e_c(x))} {\upc c V}.  \]

    \begin{itemize}
        \item Case $c = \inat$. Then our assumption says that $\vrel{\nat}{x}{V_n}$, and we need to show that
        \[ \erel {\dyn} {(\eta\, e_{\inat} (n))} {(\upc {\inat}(V_n))}. \]
        Since both are values, it suffices to show that 
        \[ \vrel{\dyn}{e_{\inat}(n)}{\upc {\inat}(V_n)}, \]
        which follows from the definition of $\vrel{\dyn}{}{}$ and our assumption.

        \item Case $c = \itimes$. Then our assumption says that $\vrel{\dyn \times \dyn}{(x, y)}{V_\times}$,
        and we need to show that
        \[ \erel {\dyn} {(\eta\, e_{\itimes}(x, y))} {(\upc {\itimes}(V_\times))}.  \]
        Since both are values, it suffices to show that
        \[ \vrel{\dyn}{e_{\itimes}(x, y)}{\upc {\itimes}(V_\times)}, \] 
        which follows from the definition of $\vrel{\dyn}{}{}$ and our assumption.
        
        % \item Case $c = \itimes(c_\times)$ where $c_\times : A_1 \times A_2 \ltdyn \dyn \times \dyn$.
        % Then our assumption says that $\vrel{A_1 \times A_2}{(x, y)}{V_\times}$, and we need to show that
        % \[ \erel {\dyn} {(\eta, (e_{\itimes(c_\times)}(x, y)))} {(\upc {\itimes(c_\times)} V_\times)}. \]
        % Since $(\upc {\itimes(c_\times)} (V_\times)) \stepsin{0} (\tagtimes (\upc c_\times (V_\times)))$,
        % then by anti-reduction it suffices to show that
        % \[ \erel {\dyn} {(\eta, (e_{\itimes(c_\times)}(x, y)))} {(\tagtimes (\upc c_\times (V_\times)))}. \]
        
        % Note that $e_{\itimes(c_\times)} = e_{\itimes} \circ e_{c_\times}$,
        % so by the induction hypothesis for $c_\times$, we have
        % \[ \erel {\dyn \times \dyn} {e_{c_\times}(x, y)} {\upc c_\times (V_\times)}. \]
        % Then the result follows using the same reasoning as in case $\itimes$.
        % Note: the terms are closed, so when we apply the IH we don't need to worry about the substitutions.
        
        \item Case $c = \iarr$.
        Then our assumption says that $\vrel{\dyn \ra \dyn}{f}{V_f}$, and we need to show that
        \[ \vrel{\dyn}{e_{\iarr}(f)}{\upc \iarr (V_f)}. \]
        By the definition of $\vrel{\dyn}{}{}$, it suffices to show that for all $x, V$ we have
        \[ (\later(\vrel{\dyn}{x}{V}) \to \erel{\dyn}{(f\, x)}{(V_f\, V)}) \]
        But we know that this holds now by our assumption, so it also holds later.

        % \item Case $c = \iarr(c_\arr)$ where $c_\arr : A_i \ra A_o \ltdyn \dyn \ra \dyn$.
        % Then our assumption says that $\vrel{A_i \ra A_o}{f}{V_f}$, and we need to show that
        % \[ \erel {\dyn} {(\eta\, (e_{\iarr(c_\arr)}(f)))} {(\upc {\iarr(c_\arr)} V_f)}. \]
        % Since $(\upc {\iarr(c_\arr)} V_f) \stepsin{0} (\tagarr (\upc c_\arr (V_f)))$, by anti-reduction it suffices to show that
        % \[ \erel {\dyn} {(\eta\, (e_{\iarr(c_\arr)}(f)))} {(\tagarr (\upc c_\arr (V_f)))}. \]
        
        % Note that $e_{\iarr(c_\arr)} = e_\tagarr \circ e_{c_\arr}$,
        % so by the induction hypothesis for $c_\arr$, we have
        % \[ \erel {\dyn \ra \dyn} {(e_{\iarr(c_\arr)}(f))} {(\upc c_\arr (V_f))}. \]
        % Then the result follows using the same reasoning as in case $\tagarr$.


        \item Case $c = c_1 c_2$ where $c_1 : A_1 \ltdyn A_2$ and $c_2 : A_2 \ltdyn A_3$.
        Then our assumption says that $\vrel{A_1}{x}{V}$, and we need to show that
        \[ \erel{A_3} {(\eta(e_{c_1 c_2}(x)))} {(\upc {c_1 c_2} V)}. \]
        Note that $e_{c_1 c_2} = e_{c_2} \circ e_{c_1}$ and that 
        $\upc {(c_1 c_2)} V \stepsin{0} (\upc c_2 (\upc c_1 V))$.
        Thus, it suffices to show that
        \[ \erel{A_3} {(\eta(e_{c_2}(e_{c_1}(x))))} {(\upc c_2 (\upc c_1 V))}. \]
        
        Observe that $(\eta(e_{c_2}(e_{c_1}(x)))) = T(e_{c_2})(\eta(e_{c_1}(x)))$.
        Thus, by the semantic bind lemma, with $\phi(\xbar) = T(e_{c_2})(\xbar)$
        and $E = \upc {c_2} \bullet$ it suffices to show that
        \[ \erel{A_2}{\eta(e_{c_1}(x))}{\upc {c_1} V}, \]
        and that for any $\vrel{A_2}{x'}{V'}$ we have % $\erel{A_3}{\phi(\eta(x'))}{E[V']}$
        \[ \erel{A_3}{\eta(e_{c_2}(x'))}{\upc {c_2} V'}. \]
        These hold by the induction hypotheses for $c_1$ and $c_2$ respectively.


        % TODO: need to justify why the value is a pair of values?
        \item Case $c = c_1 \times c_2$ where $c_1 : A_1 \ltdyn A_1'$ and $c_2 : A_2 \ltdyn A_2'$.
        Then our assumption says that $\vrel{A_1 \times A_2}{(x_1, x_2)}{V_\times}$.
        Note that since $V_\times$ is a closed value of type $A_1 \times A_2$,
        then we have $V_\times = (V_1, V_2)$, where $V_1$ and $V_2$ are closed
        values of type $A_1$ and $A_2$ respectively. We need to show that
        \[ \erel {A_1' \times A_2'} {(\eta(e_{c_1 \times c_2}(x_1, x_2)))} {(\upc {c_1 \times c_2} (V_1, V_2))}. \]
        Now note that $e_{c_1 \times c_2}(x_1, x_2) = (e_{c_1}(x_1), e_{c_2}(x_2))$.
        Furthermore, we have
        $\upc {(c_1 \times c_2)} (V_1, V_2) \stepsin{0} (\upc {c_1} V_1, \upc {c_2} V_2)$,
        so by anti-reduction, it suffices to show
        \[ \erel {A_1' \times A_2'} {(\eta(e_{c_1}(x_1), e_{c_2}(x_2)))} {(\upc {c_1} V_1, \upc {c_2} V_2)}. \]
        Now using the semantic bind lemma, with $\phi(\xbar) = T(\lambda x. (x, e_{c_2}(x_2)))(\xbar)$ and 
        $E = (\bullet, \upc c_2 V_2)$, it suffices to show that (1)
        $\erel{A_1'}{\eta(e_{c_1}(x_1))}{(\upc {c_1} V_1)}$,
        and that (2) for all $\vrel{A_1'}{x_1'}{V_1'}$, we have 
        \[ \erel{A_1' \times A_2'}{\eta(x_1', e_{c_2}(x_2))}{(V_1', \upc c_2 V_2)}. \]
        Note that (1) follows from the inductive hypothesis for $c_1$, and to show (2), we apply the
        semantic bind lemma again, with $\phi'(\xbar) = T(\lambda y. (x_1', y))$ and $E' = (V_1', \bullet)$.
        This time, it suffices to show that $\erel{A_2'}{\eta(e_{c_2}(x_2))}{(\upc c_2 V_2)}$,
        and that for all $\vrel{A_2'}{x_2'}{V_2'}$, we have
        \[ \erel{A_1' \times A_2'}{\eta(x_1', x_2')}{(V_1', V_2')}. \]
        The former follows from the inductive hypothesis for $c_2$, and the latter follows
        from the assumptions that $\vrel{A_1'}{x_1'}{V_1'}$ and $\vrel{A_2'}{x_2'}{V_2'}$. 

        \item Case $c = c_i \ra c_o$, where $c_i : A_i \ltdyn A_i'$ and $c_o :
        A_o \ltdyn A_o'$. Follows by anti-reduction, the semantic bind lemma,
        and the induction hypotheses for $c_i$ and $c_o$.
        
        % We need to show
        % %
        % \[ \erel {A_i' \ra A_o'} 
        %     {\sem{\upc {(c_i \ra c_o)} (f\, \gambar)}} 
        %     {\upc {(c_i \ra c_o)} V_f[\gamma']}.
        % \]
        % %
        % We have 
        % %
        % \begin{align*}
        %     \sem{\upc {(c_i \ra c_o)} (f\, \gambar)} 
        %         &= ((\sem{\dnc {c_i}} \tok T\sem{A_o'}) \circ (\sem{A_i} \tok UT(\sem{\upc {c_o}}))) (f\, \gambar) \\
        %         &= \lambda x'. T\sem{\upc {c_o}}[(f\, \gambar)^\dag (\sem{\dnc {c_o} (x')})],
        % \end{align*}
        % %
        % where $\tok$ is the Kleisli action of $\to$.

        % Let 
        % %
        % $\bar{M} := \lambda x'. T\sem{\upc {c_o}}[(f\, \gambar)^\dag (\sem{\dnc {c_o} (x')})]$,
        % %
        % and let
        % %
        % $N := (\lambda x. (\upc {c_o} (V_f[\gamma']\, (\dnc {c_i} x))))$.
        % %
        % Note that
        % %
        % \[ \upc {(c_i \ra c_o)} V_f[\gamma'] \stepsin{0} N. \] 
        % %
        % Thus, by the anti-reduction lemma, it suffices to show that
        % %
        % \[ \erel {A_i' \ra A_o'} {\bar{M}} {N}. \]
        % %
        
    \end{itemize}

    Now we consider the downcasts. We again apply the semantic bind lemma. Let
    $x' : \sem{A'}$ and $V' : \val{A'}$, and suppose that $\vrel{A'}{x'}{V'}$.
    We will show that
    %
    \[ \erel{A}{p_c(\eta\, x')} {\dnc c V'}.  \]

    \begin{itemize}
        \item Case $c = \inat$. Then $x' \in D$, and $V' \in \val{\dyn}$, and our
        assumption says that $\vrel{\dyn}{x'}{V'}$. We need to show that
        \[ \erel {\nat} {p_{\inat}(\eta\, x')} {\dnc {\inat} V'}.  \]
        We proceed by cases on the value $x'$ of type $D$. If it's $\inat(n)$
        for some $n$, then we know that $V' = \upc {\inat} V_n$ for some value
        $V_n$, where $V_n = \suc^n(\zro)$.
        %
        We have that $p_{\inat}(\eta\, \inat(n)) = \eta\, n$, and that $\dnc
        \inat (\upc \inat V_n) \stepsin{0} V_n$. Thus, by anti-reduction, it
        suffices to show that
        %
        \[ \erel {\nat} {\eta\, n} {V_n}. \]
        %
        It then suffices to show $\vrel{\nat}{n}{V_n}$, which holds by the fact
        that $V_n = \suc^n(\zro)$.

        If $x'$ is a product or a function, then $V'$ will be as well. Then
        $p_{\inat}(\eta\, x') = \mho$, and $\dnc \inat V' \stepsin{0} \mho$.
        Then since $\erel{\nat}{\mho}{\mho}$ by definition of the relation, this
        completes the proof in this case.

        \item Case $c = \itimes$: this is similar to the $\inat$ case.
        
        \item Case $c = \iarr$. Then $x' \in D$ and $V' \in \val{\dyn}$, and our
        assumption says that $\vrel{\dyn}{x'}{V'}$. We need to show that
        %
        \[ \erel {\dyntodyn} {p_{\iarr}(\eta\, x')} {\dnc {\iarr} V'}. \]
        %
        We proceed by cases on the value $x'$ of type $D$. If it's a natural
        number or product, then $V'$ will be as well. Then both sides will
        error, and the proof is complete.
        %
        Otherwise, if $x' = \iarr(h)$ for some $h$, then we know that $V' = \upc
        {\iarr} V_f$ for some $V_f \in \val{\dyntodyn}$, and that $\later (\vrel{\dyntodyn}{h}{V_f})$.

        We have that $p_{\iarr}(\eta\, \iarr(h)) = \theta(\nxt(\eta\, h))$, i.e., a
        step is inserted. We also know that $\dnc {\iarr} (\upc {\iarr} V_f)
        \stepsin{1} V_f$. Thus, by anti-reduction for true steps (Lemma
        \ref{lem:anti-reduction-true}), it suffices to show that
        %
        \[ \later_t(\erel{\dyntodyn}{(\nxt(\eta\, h))_t}{V_f}), \]
        i.e.,
        \[ \later (\erel{\dyntodyn}{(\eta\, h)}{V_f}). \]
        %
        This follows from the fact that $\later (\vrel{\dyntodyn}{h}{V_f})$.

        \item Case $c = c_1 c_2$: this is similar to the corresponding upcast case.
        \item Case $c = c_i \ra c_o$: this is similar to the corresponding upcast case.
        \item Case $c = c_1 \times c_2$: this is similar to the corresponding upcast case.
    \end{itemize}


\end{proof}


Now we can prove the fundamental lemma:
\begin{lemma}[Fundamental Lemma]
  Let $\Gamma \vdash M : A$. Then $\Gamma \vDash \sem{M} \sim M : A$.
\end{lemma}
\begin{proof}
    By induction on the term $M$, using the compatibility lemmas in each case.
\end{proof}



%\section{Omitted Proofs Section \ref{sec:globalization}}

\section{Omitted Details and Proofs for Section \ref{sec:framework}}

\subsection{Graduality Schemata for Free Theories}

We define by induction on execution traces satisfaction relations $M \vDash t$
between an ITree $M$ and a $\textrm{PTrace}$ $t$ and $f \vDash t$ between an
$R$-indexed family of ITrees and an $\textrm{OTrace}R$.
\begin{align*}
  {M \vDash \eta(V)} &= M \leadstomany \leaf(\eta(V))\\
  {M \vDash \mho} &= M \leadstomany \leaf(\mho)\\
  {M \vDash \sigma; o} &= M \leadstomany \sigma(f) \wedge f \vDash o\\
\end{align*}
\begin{align*}
  f \vDash \textrm{end} &= \top \\
  f \vDash r; p &= f(r) \vDash p
\end{align*}
%
Additionally, we inductively define an error ordering $\sqsubseteq$ on player
and opponent traces which intuitively means if $p \sqsubseteq q$ then $p$ is
either equal to $q$ or it is a prefix of $q$ that ends in $\mho$.
\begin{mathpar}
  %\footnotesize
  \mho \sqsubseteq p

  \eta V \sqsubseteq \eta V

  \inferrule{o \sqsubseteq o'}{\sigma; o \sqsubseteq \sigma; o'}

  {\textrm{end} \sqsubseteq \textrm{end}}

  \inferrule{p \sqsubseteq q}{r; p \sqsubseteq r; q}
\end{mathpar}

Next, we state two lemmas about the coinductive error ordering and bisimilarity
on ITrees. The proofs are by a straightforward coinductive argument.

\begin{lemma}[Properties of $\sqsubseteq$]
  \label{lem:itree-strong-err}
  Assume $M \sqsubseteq N$.
  \begin{enumerate}
  \item If $M \leadsto \eta x$ then $N \leadsto \eta x$
  \item If $M \leadsto \sigma(f)$ then $N \leadsto \sigma(g)$ with $\forall r\in \ar(\sigma). f(r) \sqsubseteq g(r)$.
  \item If $N \leadsto \err$ then $M \leadsto \err$
  \item If $N \leadsto \eta x$ then $M \leadsto \err$ or $M \leadsto \eta x$
  \item If $N \leadsto \sigma(f)$ then $M \leadsto \err$ or $M \leadsto \sigma(g)$ with $\forall r \in \ar(\sigma). f(r) \sqsubseteq g(r)$
  \end{enumerate}
\end{lemma}

\begin{lemma}[Properties of $\bisim$]
  \label{lem:itree-bisim}
  Assume $M \sqsubseteq N$.
  \begin{enumerate}
  \item If $M \leadsto \err$ then $N \leadsto \err$
  \item If $M \leadsto \eta x$ then $N \leadsto \eta x$
  \item If $M \leadsto \sigma(f)$ then $N \leadsto \sigma(g)$ with $\forall r\in \ar(\sigma). f(r) \bisim g(r)$.
  \item If $N \leadsto \err$ then $M \leadsto \err$
  \item If $N \leadsto \eta x$ then $M \leadsto \eta x$
  \item If $N \leadsto \sigma(f)$ then $M \leadsto \sigma(g)$ with $\forall r \in \ar(\sigma). f(r) \bisim g(r)$
  \end{enumerate}
\end{lemma}

Now we prove Theorem \ref{thm:itrees-ext-ord} about the extensional ordering $\ltbisim$ on ITrees:
\begin{theorem}[Properties of $\ltbisim$]
  $\ltbisim$ is transitive.
  If $M \sqsubseteq N$ then $M \ltbisim N$.
  If $M \bisim N$ then $M \ltbisim N$.
\end{theorem}
\begin{proof}
  Transitivity is clear. The other two implications follow by induction on the
  structure of the input trace and Lemmas \ref{lem:itree-strong-err} and
  \ref{lem:itree-bisim}.
\end{proof} 

\subsection{Exception Transformer for Graduality Schemata}

In this section, we show that the exception monad transformer acts on
$\calA$-graduality schemata. More concretely, the goal of this section is to
prove the following theorem:
%
\begin{theorem}[Exception Transformer for Graduality Schemata]
  Let $E$ be a finite type. Let $\calA$ be an ordinary algebraic theory, and
  suppose we are given an $\calA$-graduality schema. Then there is a graduality
  schema for the transformed theory $(\calA + \exn{E})$.
\end{theorem}

We first show that if $T$ gives the free models of $\calA$ on Set, then the
error monad transformer applied to $T$ gives the free models of the theory of
$\calA + \exn{E}$ on Set:
%
\begin{lemma}\label{lem:error-monad-transformer-free-models}
    Let $\calA$ be a guarded algebraic theory whose free models are given by
    $T$. Then the free models of the theory $\calA + \exn{E}$ are given by
    $T'(A) = T(A + E)$.
\end{lemma}
\begin{proof}
    Let $\calA = (\Sigma_\calA, \text{Eqn}_\calA)$ be a guarded algebraic theory with
    free-model monad given by $T_\calA$. Let $T'(A) = T(A + E)$.

    We first show that $T'(A)$ implements each operation in $\calA +
    \exn{E}$. For each $e : E$, we define 
    %
    \[ \sem{\ex{e}} = T.\eta(\inr\, e). \]
    %
    For each operation $\oper : (\tau, \gamma)$ in $A$, we define $\sem{\oper} :
    ((\tau_\oper \to T'(A)) \times (\gamma_\oper \to \later (T'(A)))) \to T'(A)$
    using the fact that $T(A+E)$ implements the operations, i.e, we define
    %
    \[ \sem{\oper}(u, g) = T(A+E).\sem{\oper}(u, g) \]
    
    We define the inclusion of generators $\eta' : A \to T'(A)$ by
    %
    \[ \eta'(x) = T(A+E).\eta(\inl\, x). \]

    Next, we show that the equations are satisfied. Since there are no equations
    in $\exn{E}$, we need only show that the equations of $\calA$ hold.
    Specifically, let $e = (\Gamma, l, r)$ be an equation in $\text{Eqn}_\calA$, where
    $l, r$ are terms in the signature $\Sigma_\calA$ with variables in $\Gamma$.
    Consider an arbitrary $f : \Gamma \to T'(A)$. We need to show that
    $\sem{l}_f = \sem{r}_f$. Since we know that $l$ and $r$ do not mention the operations
    $\ex{e}$, we know that this equality holds by the fact that $T(A + E)$ is a
    model for $\calA$.

    Finally, we show the universal property. Let $B$ be a model for $\calA +
    \exn{E}$ and let $f : A \to |B|$. We define $\fbar : T'(A) \to |B|$ using
    the universal property of $T(A + E)$. Namely, since $B$ is also a model for
    $\calA$, it suffices to define a function $g : A + E \to |B|$. We define
    %
    \[ g(\inl\, x) = f(x), \] 
    and
    \[ g(\inr\, e) = B.\ex{e}. \]
    %
    Then by the universal property this extends uniquely to a function $\ghat :
    T(A + E) \to |B|$ that preserves the operations in $\calA$.
    %
    It is clear that this is a homomorphism of models of $\calA +
    \exn{E}$ since for each $e : E$, we have $g(\sem{\ex{e}}) = \ghat(T.\eta(\inr\, e)) = g(\inr\, e) =
    B.\ex{e}$.

    For uniqueness, suppose $h, h' : T(A + E) \to B$ are homomorphisms of
    $(\calA + \exn{E})$-models such that $h \circ \eta' = h' \circ \eta'$. We
    show that $h = h'$. Since $h$ and $h'$ are in particular homomorphisms of
    $\calA$-models, by the universal property for $T(A + E)$, it suffices to
    show that they agree when composed with $T(A+E).\eta$. Thus, let $x \in A +
    E$. We proceed by cases on $x$. If $x = \inl\, a$, then we have
    %
    \[ 
      h (\eta (\inl\, a)) 
      = h (\eta'(a)) 
      = h' (\eta'(a)) 
      = h' (\eta (\inl\, a)). 
    \]
    %
    Now suppose $x = \inr\, e$. Then we have
    \[ 
      h (\eta (\inr\, e)) 
      = h(\sem{\ex{e}}) 
      = B.\ex{e} 
      = h'(\sem{\ex{e}}) 
      = h' (\eta (\inr\, e)).  
    \]

\end{proof}

Now we show that the monad transformer lifts to relations:
%
\begin{lemma}\label{lem:error-monad-transformer-free-domains}
    Let $\calA$ be an ordinary algebraic theory. Let $\calA^G \in \calA[\mho,
    \theta]$, i.e., $\calA^G$ is a guarded extension of $\calA$ with an error
    and a $\theta$ operation. Suppose that the adjunction between sets and
    models of $\calA^G$ lifts to a strong adjunction between predomains and
    $\calA^G$-domains.
    
    Then for the theory $\calA^{G+} := (\calA^{G} + \exn{E})$, the adjunction
    between sets and models of $\calA^{G+}$ lifts to a strong adjunction between
    predomains and $\calA^{G+}$-domains.
\end{lemma}
\begin{proof}
    % We define $\calA^{G+}$ to be $\calA^{G} + \mho$. % i.e., exceptions
    Let the free models of $\calA^G$ be given by $T$, and let the free domains
    of $\calA^G$ be given by $\That$. By assumption, for any predomain $A$, we
    have $|\That(A)| = T(|A|)$.
    
    We know by Lemma \ref{lem:error-monad-transformer-free-models} that at the
    level of sets, the free models of $\calA^{G+}$ are given by $T'(A) = T(A +
    E)$.
    
    We show that $\That'(A) := \That(A + E)$ gives the free domains of
    $\calA^{G+}$. In $A + E$, the $+$ symbol denotes the coproduct of
    predomains, where in the ordering and bisimilarity relations, $\inl$ is
    related to $\inl$ and $\inr$ to $\inr$, and since $E$ is merely a type, the
    ordering and bisimilarity are equality.

    The proof proceeds in several steps:
    \begin{itemize}
        \item We define the strong error ordering and weak bisimilarity on
        $\That'(A)$ to be those of $\That(A + E)$. These relations satisfy the
        properties required of the strong error ordering and weak bisimilarity
        relations, since the relations on $\That(A)$ satisfy the properties for
        all $A$.

        \item For a predomain relation $c : A \rel A'$, we define the lifting
        $\That'(c) = \That(c + =_E)$.

        \item We check the universal property of the lifting, as follows.
        Consider a square $f \ltsq{c}{Ud} g$ where $c : A \rel A'$ is a
        predomain relation, and $d : B \rel B'$ is a relation of
        $\calA^{G+}$-domains.
        %
        Then we need to show that $f^* \ltsq{\That'(c)}{d} g^*$, where $-^*$
        denotes the extension operation turning a function $A \to UB$ into a
        homomorphism $T'(A) \to B$. By the definition of $-^*$ and the universal
        property of the lifting for $\That$, it suffices to construct a
        predomain square
        %
        $(f + B.\ex{\cdot}) \ltsq{c + =_E}{UB} (g + B.\ex{\cdot})$.
        %
        We proceed by cases on the input. In case both are $\inl$, then we use
        the square $f \ltsq{c}{Ud} g$. If both are $\inr e$, then we need to
        show that $(B.\ex{e}) \binrel{d} (B'.\ex{e})$ which is true as $d$ is a
        relation of $\calA^{G+}$ domains and hence is a congruence with respect
        to $\ex{e}$.


        \item It is easily checked that the operations are monotone and
        bisimilarity-preserving, and that the strength
        $t : A \times T'(A') \to T'(A \times A')$ is a morphism of predomains.
        \item Lastly, for any predomain $A$, we have   
        \begin{align*}
            |\That'(A)| 
            &= |\That(A + E)| \\
            &= T(|A + E|)\\ 
            &= T(|A| + E)\\ 
            &= T'(|A|).
        \end{align*}
    \end{itemize}
\end{proof}

Now we are ready to define the action of the exception transformer on graduality
schemata. Let $\calA$ be an ordinary algebraic theory, and let $\calG$ be a
graduality schema for $\calA$.
%
We write $\calAs, \calAsb$ for the guarded extensions of $\calA$ associated with
$\calG$, and we write $T$ for the free-model monad for $\calAs$ on Set, and
$\That$ for its lifting to predomains.
%
We define a graduality schema for $(\calA + \exn{E})$, as follows.
\begin{itemize}
    \item We define the guarded extensions 
    $\calAgplus := \calAs + \exn{E}$ and 
    $\calAgplus_b := \calAsb + \exn{E}$.

    \item Then by Lemma \ref{lem:error-monad-transformer-free-models}, the free models of
    $\calAgplus$ on Set are given by $T'(A) := T(A + E)$.

    \item By Lemma \ref{lem:error-monad-transformer-free-domains}, we have that 
    the free domains for $\calAgplus$ are given by $\That(A + E)$.

    \item We define the extensional types $Z'(A) := Z(A + E)$ where $Z$ is the
    family of extensional types for $\calG$ (we use $Z$ here to avoid confusion
    with the type $E$ of exceptions). Next, we define the function
    %
    $\col'_A : T'^{gl}(A) \to Z'(A)$, i.e., $T^{gl}(A + E) \to Z(A + E)$ as 
    %
    \[ \col'_A := \col_{A + E}, \]
    %
    where $\col$ is the family of functions associated with the schema $\calG$.
    
    \item Lastly, the relation $\ltbisim_{Z'(A)}$ is defined to be
    $\ltbisim_{Z(A + E)}$. Then since each $\ltbisim_{Z(A)}$ is a partial order,
    in particular $\ltbisim_{Z(A + E)}$ is a partial order. Furthermore, it is
    clear that the error ordering and bisimilarity relations on the globalized
    monad imply $\ltbisim_{Z'(A)}$ on the output of $\col'_A$, since this holds for
    $\ltbisim_{Z(A)}$ for all $A$.
\end{itemize}



\subsection{State Transformer for Graduality Schemata}

In this section, we show that the state monad transformer acts on
$\calA$-graduality schemata. We prove the following theorem:
\begin{theorem}[Global State Transformer for Graduality Schemata]
  Let $S$ be a finite type. Let $\calA$ be an ordinary algebraic theory, and
  suppose we are given an $\calA$-graduality schema. Then there is a graduality
  schema for the transformed theory $(\calA \otimes \state{S})$.
\end{theorem}

% In particular:
% \begin{itemize}
%     \item For any guarded algebraic theory $\calA$, if the free models of
%     $\calA$ are given by $T$, then the free models of $\calA \otimes \state{S}$
%     are given by $S \to T(S \times A)$.
% \end{itemize}

First, we show that for any guarded algebraic theory $\calAg$, if the free
models of $\calAg$ are given by $T$, then the free models of $\calAg \otimes
\state{S}$ are given by $S \to T(S \times A)$. (This is Theorem
\ref{thm:state-monad-transformer} in the paper.)

Let $\calAg = (\Sigma, E)$ be a guarded algebraic theory with free-model monad
$T$. We will show that the free-model monad for the theory $\calAg \otimes
\state{S}$ is given by
%
\[ T'(A) := S \to T(S \times A). \]
%
We begin by defining the operations. We let

\[ \sem{\get} f\, s = f\, s\, s, \]
%
\[ \sem{\putop_s} m\, s' = m\, s, \]  
%
and we define $\sem{\oper} : [(\tau_\oper \to T'(A)) \times (\gamma_\oper \to \later(T'(A)))] \to T'(A)$ by
%
\[ 
  \sem{\oper}(u, g)\, s = 
    T(A \times S).\sem{\oper}
      [(\lambda w. u\, w\, s), (\lambda z. \lambda t. g\, z\, t\, s)].
\]

Now, we show that the required equations hold when interpreted in $T'(A)$.
%
The proofs of the get-put laws are analogous to those for the ordinary state
monad and hence omitted.
%
We prove the tensor equations stating that $\get$ and $\putop$ commute with each
operation in $\calA$:
\begin{enumerate}
    \item \emph{Get commutes with each operation in $\calAg$}: Let $\oper :
    (\tau, \gamma) \in \calAg$. Let $u : \tau \to S \to T'(A)$ and $g : \gamma \to
    \later(S \to T'(A))$. Let
    %
    \[ 
        M = \sem{\oper}[
              (\lambda w. \sem{\get}(\lambda s. u\, w\, s)),
              (\lambda z. \lambda t.\sem{\get}(\lambda s. g\, z\, t\, s))]
    \] 
    and
    \[
        N = \sem{\get}[\lambda s. 
              \sem{\oper}((\lambda w. u\, w\, s), (\lambda z. \lambda t. g\, z\, t\, s))]
    \]
    %
    We need to show that $M = N$. For brevity, let
    %
    $u_1 = (\lambda w. \sem{\get}(\lambda s. u\, w\, s))$ and let 
    %
    $g_1 = (\lambda z. \lambda t.\sem{\get}(\lambda s. g\, z\, t\, s))$.
    %
    We first observe that for any $w : \tau$ and $z : \gamma$, we have
    %
    $u_1(w) = (\lambda s. u\, w\, s\, s)$ of type $S \to T(S \times A)$ and 
    %
    $g_1(z) = (\lambda t. \lambda s. g\, z\, t\, s\, s)$ of type $\later (S \to T(S \times A))$.
    
    Now let $s \in S$. We have
    \begin{align*}
        M(s) &= \sem{\oper}(u_1, g_1)(s) \\
             &= T(A \times S).\sem{\oper}[(\lambda w. u_1\, w\, s), (\lambda z. \lambda t. g_1\, z\, t\, s)] \\
             &= T(A \times S).\sem{\oper}[(\lambda w. u\, w\, s\, s), (\lambda z. \lambda t. g\, z\, t\, s\, s) ]. \\
    \end{align*}
    
    Now let $p(s) = (\lambda w. u\, w\, s)$ and $q(s) = (\lambda z. \lambda t. g\, z\, t\, s)$.
    We have
    %
    \begin{align*}
        N(s) &= (\sem{\get}[\lambda s'. \sem{\oper}(p(s'), q(s'))])(s) \\
             &= (\lambda s'. \sem{\oper}(p(s'), q(s')))(s)(s) \\
             &= [\sem{\oper}(p(s), q(s))](s) \\
             &= T(A \times S).\sem{\oper}[(\lambda w. p(s)\, w\, s), (\lambda z. \lambda t. q(s)\, z\, t\, s)] \\
             &= T(A \times S).\sem{\oper}[(\lambda w. u\, w\, s\, s), (\lambda z. \lambda t. g\, z\, t\, s\, s)].
    \end{align*}
    %
    Thus, we conclude that $M = N$.
    
    \item \emph{Put commutes with each operation in $\calAg$}: This is similar to
    the above and hence omitted.
\end{enumerate}

Finally, we prove that the equations in $\calAg$ hold. For this, we first prove
the following lemma:
%
\begin{lemma}
    % Let $t$ be a term in the signature $F_\Sigma(A)$ (recall $\Sigma$ is the
    % signature of $\calA$), and let $f : \Gamma \to T'(A)$ be a function. 

    Let $f : \Gamma \to T'(A)$ be a function. Then for all $s \in S$, there exists a
    function $\fhat_s : \Gamma \to T(S \times A)$ such that, for all terms $t$ in
    $F_\Sigma(\Gamma)$, we have
    %
    \[ \sem{t}_{\fhat_s}^{T(S \times A)} = (\sem{t}_f^{T'(A)})(s). \]
\end{lemma}

\begin{proof}
    For $s \in S$, we define $\fhat_s(x) = f(x)(s)$.
    %
    The equality follows by induction on the term $t$:
    \begin{itemize}
        \item If $t$ is a variable $x \in \Gamma$, then we have
        %
        \[ \sem{x}_{\fhat_s}^{T(S \times A)} 
          = \fhat_s(x) 
          = f(x)(s) 
          = (\sem{x}_f^{T'(A)})(s).
        \]

        \item Now suppose $t = \oper(u, g)$, where (1) $\oper \in \Sigma$, 
        (2) $u : \tau_\oper \to F_\Sigma(\Gamma)$ and 
        (3) $g : \gamma_\oper \to \later (F_\Sigma(\Gamma))$. Then the left-hand side is
        %
        \begin{align*}
            \sem{\oper(u, g)}_{\fhat_s}^{T(S \times A)}
              &= T(S \times A).\sem{\oper}(u_1, g_1),
        \end{align*}
        %
        where $u_1(w) := \sem{u(w)}_{\fhat_s}^{T(S \times A)}$ and 
        $g_1(z)(t) := \sem{g(z)(t)}_{\fhat_s}^{T(S \times A)}$.

        The right-hand side is
        \begin{align*}
            (\sem{\oper(u, g)}_f^{T'(A)})(s)
              &= [T'(A).\sem{\oper}(u_2, g_2)](s) \\
              &= T(S \times A).\sem{\oper}[
                  (\lambda w. u_2\, w\, s), 
                  (\lambda z. \lambda t. g_2\, z\, t\, s)],
        \end{align*}
        where $u_2(w) := \sem{u(w)}_f^{T'(A)}$ and $g_2(z)(t) := \sem{g(z)(t)}_f^{T'(A)}$.
        
        Then by the inductive hypothesis, we have $u_1(w) = u_2(w)(s)$ and
        $g_1(z)(t) = g_2(z)(t)(s)$, so the above becomes
        %
        \[ T(S \times A).\sem{\oper}[
                  (\lambda w. u_1\, w), 
                  (\lambda z. \lambda t. g_1\, z\, t)].
        \]
        %
        Thus, both sides are equal.

    \end{itemize}
\end{proof}

Now we can show that the equations in $\calAg$ hold:
\begin{proof}
    Let $e = (\Gamma, l, r) \in
    E$, and let $f : \Gamma \to T'(A)$ be a function. We need to show

    \[ \sem{l}_f^{T'(A)} = \sem{r}_f^{T'(A)}, \]

    where both sides are functions $S \to T(S \times A)$. To this end, let $s \in
    S$. Then by the above Lemma, we have that there is $\fhat_s : \Gamma \to T(S \times
    A)$ with the property described in the Lemma. In particular, we have
    %
    \begin{align*}
        (\sem{l}_f^{T'(A)})(s) 
        &= \sem{l}_{\fhat_s}^{T(S \times A)} \\
        &= \sem{r}_{\fhat_s}^{T(S \times A)} \\
        &= (\sem{t}_f^{T'(A)})(s),
    \end{align*}
    %
    where the second equality holds because $T(S \times A)$ is a model of $\calAg$
    and hence the equation holds when interpreted in $T(S \times A)$.
\end{proof}

Next, we define $\eta' : A \to T'(A)$ by 
%
$\eta'(x) = \lambda s.\eta(s, x)$, 
%
where $\eta : (S \times A) \to T(S \times A)$ is the unit for $T$.

Finally, we show the universal property.
\begin{proof}
    Let $B$ be a model of the theory $\calAg \otimes \state{S}$, and let $f : A \to UB$.
    %
    First, define $g : A \times S \to B$ as the following composition:
    %
    \[ g = A \times S \xrightarrow{f \times \id} UB \times S \xrightarrow{B.\sem{\putop} B} \]
    %
    We then define $\fbar : T'(A) \to B$ as the following composition:
    %
    \[ \fbar = 
        S \to T(S \times A) \xrightarrow{S \to \gbar} S \to B \xrightarrow{B.\sem{\get}} B.  \]
    %
    where $\gbar$ is the unique extension of $g$ to a homomorphism of $\calAg$-models.


    Then it is easily checked that $\fbar$ is a homomorphism of models of $\calAg
    \otimes {\state{S}}$, and that $\fbar \circ \eta' = f$.
\end{proof}

Now we show uniqueness. Let $\phi, \phi' : T'(A) \to B$ be homomorphisms of
$(\calAg \otimes {\state{S}})$-models, and suppose $\phi(\eta'\, x) =
\phi'(\eta'\, x)$ for all $x$. That is, $\phi(\lambda s. \eta(s,x)) =
\phi'(\eta(s,x))$. Then we claim that $\phi = \phi'$.

First, we need a small lemma:
\begin{lemma}
    For all $m : T(S \times A)$, we have $\phi(\lambda s.m) = \phi'(\lambda s.m)$.
\end{lemma}
\begin{proof}
    Define $\psi, \psi' : T(S \times A) \to B$ by $\psi(m) = \phi(\lambda s.m)$
    and $\psi'(m) = \phi'(\lambda s.m)$. First, note that $\psi$ and $\psi'$ are
    homomorphisms of $\calAg$-models.
    %
    Thus, by the universal property of $T(S \times A)$ as the free-model monad
    for $\calAg$, it suffices to show that $\psi \circ \eta = \psi' \circ \eta$.
    We have
    \begin{align*}
        \psi(\eta(s,a)) 
        &= \phi(\lambda s'.\eta(s, a)) \\
        &= \phi(\putop_s (\lambda s''.\eta(s'', a))) \\
        &= B.\putop_s(\phi(\lambda s''.\eta(s'', a))) \\
        &= B.\putop_s(\phi'(\lambda s''.\eta(s'', a))) \\
        &= \phi'(\putop_s(\lambda s''.\eta(s'', a))) \\
        &= \phi'(\lambda s'.\eta(s, a)) \\
        &= \psi'(\eta(s,a)).
    \end{align*}
\end{proof}

Now let $x : T'(A)$. We have
\begin{align*}
    \phi(x)
    &= \phi(\get(\lambda s. \putop_s(x))) \\
    &= B.\get(\lambda s. \phi(\putop_s(x))) \\
    &= B.\get(\lambda s. \phi(\lambda s'. x\, s)) \\
    &= B.\get(\lambda s. \phi'(\lambda s'. x\, s)) \\
    &= B.\get(\lambda s. \phi'(\putop_s(x))) \\
    &= \phi'(\get(\lambda s.\putop_s(x))) \\
    &= \phi'(x).
\end{align*}

Next, we show that the monad transformer lifts to relations:
\begin{lemma}\label{lem:state-monad-transformer-free-domains} Let $\calA$ be an
    ordinary algebraic theory. Let $\calA^G \in \calA[\mho, \theta]$, be a
    guarded extension of $\calA$. Suppose that the adjunction between sets and
    models of $\calA^G$ lifts to a strong adjunction between predomains and
    $\calA^G$-domains, with the free domains being given by $\That$.
    
    Then for the theory $\calA^{G\otimes} := (\calA^{G} \otimes \state{S})$, the
    adjunction between sets and models of $\calA^{G\otimes}$ lifts to a strong
    adjunction between predomains and $\calA^{G\otimes}$-domains.
    The free domains for $\calA^{G\otimes}$ are given by $S \to \That(S \times A)$.
\end{lemma}
\begin{proof}
    % We define $\calA^{G+}$ to be $\calA^{G} + \mho$. % i.e., exceptions
    Let the free models of $\calA^G$ be given by $T$, and let the free domains
    of $\calA^G$ be given by $\That$. By assumption, for any predomain $A$, we
    have $|\That(A)| = T(|A|)$.
    
    We know by the above that at the level of sets, the free models of
    $\calA^{G+}$ are given by $T'(A) = S \to T(S \times A)$.
    
    We show that $\That'(A) := S \to \That(S \times A)$ gives the free domains
    of $\calA^{G+}$. Here, $S$ is a discrete predomain, so the relations are
    given by equality. In addition, $\to$ denotes the predomain of morphisms of
    predomains between $S$ and $\That(S \times A)$ (this is the same as ordinary
    functions from $S$ to $\That(S \times A)$), and $\times$ denotes the product
    of predomains.

    \begin{itemize}
        \item The strong error ordering and weak bisimilarity on $\That'(A)$ are
        those of the predomain $S \to \That(S \times A)$.

        \item For a predomain relation $c : A \rel A'$, the lifting is given by
        $\That'(c) = S \to \That(S \times c)$, where $\That(S \times c)$ is the
        lifting of the relation $=_S \times c$.
        %
        More concretely, we have $f \binrel{\That'(c)} g$ iff for all $s \in S$,
        we have $(f\, s) \binrel{\That(S \times c)} (g\, s)$.

        \item To verify the universal property of the lifting of a relation, we
        use the definition of the unique extension of a function $A \to UB$ to a
        homomorphism out of $T'(A)$, as defined in the proof of Theorem
        \ref{thm:state-monad-transformer} above. The universal property for the
        lifted relation then follows using the actions of $\to$ and $\times$ on
        squares as well as the universal property of $\That$ on relations.

        \item It is easily checked that the operations are monotone and
        bisimilarity-preserving, and that the strength $t : A \times T'(A') \to
        T'(A \times A')$ is a morphism of predomains.
        
        \item Lastly, for any predomain $A$, we have
        %
        \begin{align*}
            |\That'(A)| &= |S \to \That(S \times A)| \\
                        &= S \to |\That(S \times A)| \\
                        &= S \to T(|S \times A|)  \\
                        &= T'(|A|).
        \end{align*}
    \end{itemize}
\end{proof}

Now we are ready to define the action of the global state transformer on
graduality schemata. Let $\calA$ be an ordinary algebraic theory, and let
$\calG$ be a graduality schema for $\calA$.
%
We write $\calAs, \calAsb$ for the guarded extensions of $\calA$ associated with
$\calG$, and we write $T$ for the free-model monad for $\calAs$ on Set, and
$\That$ for its lifting to predomains.
%
We define a graduality schema for $(\calA \otimes \state{S})$, as follows.
\begin{itemize}
    \item We define the guarded extensions 
    $\calAgtimes := \calAs \otimes \state{S}$ and 
    $\calAgtimes_b := \calAsb \otimes \state{S}$.

    \item Then by Theorem \ref{thm:state-monad-transformer}, the free models of
    $\calAgtimes$ on Set are given by $T'(A) := S \to T(S \times A)$.

    \item By Lemma \ref{lem:state-monad-transformer-free-domains}, we have that 
    the free domains for $\calAgtimes$ are given by $S \to \That(S \times A)$.

    % \item The functor $T'$ commutes with clock quantification. First note that
    % $S$ is finite and is hence clock-irrelevant. Then it follows from Theorem
    % 4.2 of kristensen et al. \cite{kristensen-mogelberg-vezzosi2022}, that $T'$
    % commutes with clock quantification.
    % %
    % Then if $A$ is clock-irrelevant, we have $\forall k. T'_k(A) \cong S \to T'$ provided $A$ is clock-irrelevant.

    \item We define the extensional types $E'(A) := S \to E(S \times A)$, where
    $E$ is the family of extensional types for $\calG$. Next, we define the
    family of functions
    %
    $\col'_A : T'^{gl}(A) \to E'(A)$, i.e., $(S \to T^{gl}(S \times A)) \to (S \to E(S \times A))$, via 
    %
    \[ \col'_A(f) := \col_{S \times A} \circ f, \]
    %
    where $\col$ is the family of functions associated with the schema $\calG$.

    \item Lastly, the relation $\ltbisim_{E'(A)}$ is defined to be $S \to
    \ltbisim_{E(S \times A)}$, i.e., we have $f \binrel{\ltbisim_{E'(A)}} g$ iff
    for all $s \in S$, we have $(f\, s) \binrel{\ltbisim_{E(S \times A)}} (g\, s)$.
    
    Then we observe that $\ltbisim_{E'(A)}$ is a partial order for each $A$. We
    now show that the error ordering and bisimilarity relations on the
    globalized monad imply $\ltbisim_{E'(A)}$ on the output of $\col'_A$.
    %
    More concretely, let $f, g : S \to T^{gl}(S \times A)$, and suppose that 
    $f \ltdyn_{S \to T^{gl}(S \times A)} g$. Then we want to show that
    $\col'_A(f) \ltbisim_{S \to E(S \times A)} \col'_A(g)$. Let $s \in S$.
    It suffices to show
    %
    \[ \col_{S \times A}(f(s)) \ltbisim_{E(S \times A)} \col_{S \times A}(g(s)). \] 
    %
    This follows from the fact that $f(s) \ltdyn_{T^{gl}(S \times A)} g(s)$.

    The same reasoning holds for bisimilarity.
\end{itemize}

\subsection{Writer Transformer for Graduality Schemata}

In this section, we show that the writer monad transformer acts on
$\calA$-graduality schemata. We prove the following theorem:
\begin{theorem}[Writer State Transformer for Graduality Schemata]
  Let $W$ be a monoid. Let $\calA$ be an ordinary algebraic theory, and
  suppose we are given an $\calA$-graduality schema. Then there is a graduality
  schema for the transformed theory $(\calA \otimes \writer{W})$.
\end{theorem}

We first show that if $T$ gives the free models of $\calA$, then the writer monad
transformer applied to $T$ gives the free models of the theory of $\calA \otimes
\writer{W}$:
\begin{lemma}\label{lem:writer-monad-transformer-free-models}
    Let $\calA$ be a guarded algebraic theory whose free models are given by
    $T$. Then the free models of the theory $\calA \otimes \writer{W}$ are given by
    $T'(A) = T(W \times A)$.
\end{lemma}
\begin{proof}
    Let $\calA = (\Sigma_\calA, \text{Eqn}_\calA)$ be a guarded algebraic theory with
    free-model monad given by $T_\calA$. Let $T'(A) = T(W \times A)$.

    We first show that $T'(A)$ implements each operation in $\calA \otimes
    \writer{W}$. For each $w : W$, we define 
    %
    \[ \sem{\wrt{w}} = T(\lambda(w',a).\, (w'w, a)). \]
    %
    For each operation $\oper : (\tau, \gamma)$ in $A$, we define $\sem{\oper} :
    ((\tau_\oper \to T'(A)) \times (\gamma_\oper \to \later (T'(A)))) \to T'(A)$
    using the fact that $T(W \times A)$ implements the operations, i.e, we define
    %
    \[ \sem{\oper}(u, g) = T(W \times A).\sem{\oper}(u, g), \]
    %
    where $T$ is the functorial action of the monad.
    
    We define the inclusion of generators $\eta' : A \to T'(A)$ by
    %
    \[ \eta'(x) = T(W \times A).\eta(\epsilon, x), \]
    %
    where $\epsilon$ is the unit of the monoid $W$.

    Next, we show that the equations are satisfied. First, the writer equations
    $\wrt{\epsilon}(x) = x$ and $\wrt{w_1}(\wrt{w_2}(x)) = \wrt{w_1 \cdot
    w_2}(x)$ follow by functoriality and associativity of the monoid.

    Now we show that the equations of $\calA$ hold. Specifically, let $e =
    (\Gamma, l, r)$ be an equation in $\text{Eqn}_\calA$, where $l, r$ are terms
    in the signature $\Sigma_\calA$ with variables in $\Gamma$. Consider an
    arbitrary $f : \Gamma \to T'(A)$. We need to show that $\sem{l}_f =
    \sem{r}_f$. Since we know that $l$ and $r$ do not mention the operations
    $\wrt{w}$, we know that this equality holds by the fact that $T(W \times A)$
    is a model for $\calA$.

    Lastly, we claim that each operation $\wrt{w}$ commutes with each operation
    $\oper$ in $\calA$. This follows from the fact that the morphism $T(W \times
    A).\sem{\oper}(u, g)$ is a homomorphism, i.e., preserves the operations in
    $\calA$.

    Finally, we show the universal property. Let $B$ be a model for $\calA \otimes
    \writer{W}$ and let $f : A \to |B|$. We define $\fbar : T'(A) \to |B|$ using
    the universal property of $T(W \times A)$. Namely, since $B$ is also a model for
    $\calA$, it suffices to define a function $g : W \times A \to |B|$. We define
    %
    \[ g(w,x) = B.\wrt{w}(f(x)). \] 
    %
    Then by the universal property this extends uniquely to a function $\ghat :
    T(W \times A) \to |B|$ that preserves the operations in $\calA$.
    %
    Then using the universal property of $T(W \times A)$ and the fact that $B$
    satisfies the associativity equation of $\writer{W}$, we can show that
    $\ghat$ is a homomorphism of models of $\calA \otimes \writer{W}$.
    %
    It is also straightforward to show that $\ghat \circ \eta' = f$, using the
    fact that $B$ satisfies the identity equation of $\writer{W}$.

    Uniqueness of $\ghat$ follows from the corresponding property for $T$.
\end{proof}


Next, we show that the monad transformer lifts to relations:
\begin{lemma}\label{lem:writer-monad-transformer-free-domains}
    Let $\calA$ be an ordinary algebraic theory. Let $\calA^G \in \calA[\mho,
    \theta]$, be a guarded extension of $\calA$. Suppose that the adjunction
    between sets and models of $\calA^G$ lifts to a strong adjunction between
    predomains and $\calA^G$-domains, with the free domains being given by
    $\That$.
    
    Then for the theory $\calA^{G\otimes} := (\calA^{G} \otimes \writer{W})$, the
    adjunction between sets and models of $\calA^{G\otimes}$ lifts to a strong
    adjunction between predomains and $\calA^{G\otimes}$-domains.
    The free domains for $\calA^{G\otimes}$ are given by $\That(W \times A)$.
\end{lemma}
\begin{proof}
    Let the free models of $\calA^G$ be given by $T$, and let the free domains
    of $\calA^G$ be given by $\That$. By assumption, for any predomain $A$, we
    have $|\That(A)| = T(|A|)$.
    
    We know by the above that at the level of sets, the free models of
    $\calA^{G+}$ are given by $T'(A) = T(W \times A)$.
    
    We show that $\That'(A) := \That(W \times A)$ gives the free domains of
    $\calA^{G+}$. Here, $W$ is a discrete predomain, so the relations are given
    by equality. In addition, $\times$ denotes the product of predomains.

    \begin{itemize}
        \item The strong error ordering and weak bisimilarity on $\That'(A)$ are
        those of the predomain $\That(W \times A)$.

        \item For a predomain relation $c : A \rel A'$, the lifting is given by
        $\That'(c) = \That(W \times c)$, i.e., the lifting of the relation $=_W \times c$.

        \item To verify the universal property of the lifting of a relation, we
        use the definition of the unique extension of a function $A \to UB$ to a
        homomorphism out of $T'(A)$, as defined in the proof of Lemma
        \ref{lem:writer-monad-transformer-free-models} above. The universal
        property for the lifted relation then follows using the action of
        $\times$ on squares as well as the universal property of $\That$ on
        relations.

        \item It is easily checked that the operations are monotone and
        bisimilarity-preserving, and that the strength $t : A \times T'(A') \to
        T'(A \times A')$ is a morphism of predomains.
        
        \item Lastly, for any predomain $A$, we have
        %
        \begin{align*}
            |\That'(A)| &= |\That(W \times A)| \\
                        &= T(|W \times A|) \\
                        &= T(W \times |A|)  \\
                        &= T'(|A|).
        \end{align*}
    \end{itemize}
\end{proof}

Now we are ready to define the action of the writer transformer on
graduality schemata. Let $\calA$ be an ordinary algebraic theory, and let
$\calG$ be a graduality schema for $\calA$.
%
We write $\calAs, \calAsb$ for the guarded extensions of $\calA$ associated with
$\calG$, and we write $T$ for the free-model monad for $\calAs$ on Set, and
$\That$ for its lifting to predomains.
%
We define a graduality schema for $(\calA \otimes \writer{W})$, as follows.
\begin{itemize}
    \item We define the guarded extensions 
    $\calAgtimes := \calAs \otimes \writer{W}$ and 
    $\calAgtimes_b := \calAsb \otimes \writer{W}$.

    \item Then by Lemma \ref{lem:writer-monad-transformer-free-models}, the free models of
    $\calAgtimes$ on Set are given by $T'(A) := S \to T(S \times A)$.

    \item By Lemma \ref{lem:writer-monad-transformer-free-domains}, we have that 
    the free domains for $\calAgtimes$ are given by $\That(W \times A)$.

    \item We define the extensional types $E'(A) := E(W \times A)$, where
    $E$ is the family of extensional types for $\calG$. Next, we define the
    family of functions
    %
    $\col'_A : T'^{gl}(A) \to E'(A)$, i.e., $(T^{gl}(W \times A)) \to (E(W \times A))$, via 
    %
    \[ \col'_A := \col_{W \times A} \]
    %
    where $\col$ is the family of functions associated with the schema $\calG$.

    \item Lastly, the relation $\ltbisim_{E'(A)}$ is defined to be $\ltbisim_{E(W \times A)}$.
    %
    Then we observe that $\ltbisim_{E'(A)}$ is a partial order for each $A$. We
    now show that the error ordering and bisimilarity relations on the
    globalized monad imply $\ltbisim_{E'(A)}$ on the output of $\col'_A$.
    %
    More concretely, let $x, y : T^{gl}(W \times A)$, and suppose that 
    $x \ltdyn_{T^{gl}(W \times A)} y$. Then we want to show that
    $\col'_A(x) \ltbisim_{E(W \times A)} \col'_A(y)$.
    But by definition of $\col'$, this is the same as showing
    %
    \[ \col_{W \times A}(x) \ltbisim_{E(W \times A)} \col_{W \times A}(y). \] 
    %
    This follows from the fact that $x \ltdyn_{T^{gl}(W \times A)} y$.

    The same reasoning holds for bisimilarity.
\end{itemize}


\subsection{Reader Transformer for Graduality Schemata}

In this section, we show that the reader monad transformer acts on
$\calA$-graduality schemata. We prove the following theorem:
\begin{theorem}[Reader Transformer for Graduality Schemata]
  Let $R$ be a finite type. Let $\calA$ be an ordinary algebraic theory, and
  suppose we are given an $\calA$-graduality schema. Then there is a graduality
  schema for the transformed theory $(\calA \otimes \reader{R})$.
\end{theorem}

First, we relate the free models of a guarded algebraic theory $\calAg$ to those
of $\calAg \times \reader{R}$.
\begin{lemma}\label{lem:reader-monad-transformer-free-models} Let $\calA$ be a
    guarded algebraic theory whose free models are given by $T$. Then the free
    models of the theory $\calA \otimes \reader{R}$ are given by $T'(A) = R \to
    T(A)$.
\end{lemma}
\begin{proof}

Let $\calAg = (\Sigma, E)$ be a guarded algebraic theory with free-model monad
$T$. We will show that the free-model monad for the theory $\calAg \otimes
\state{S}$ is given by
%
\[ T'(A) := R \to T(A). \]
%
We begin by defining the operations. We let
%
$\sem{\lookup} f\, r = f\, r\, r,$
%
and we define $\sem{\oper} : [(\tau_\oper \to T'(A)) \times (\gamma_\oper \to \later(T'(A)))] \to T'(A)$ by
%
\[ 
  \sem{\oper}(u, g)\, r = 
    T(A).\sem{\oper}
      [(\lambda w. u\, w\, r), (\lambda z. \lambda t. g\, z\, t\, r)].
\]

Now, we show that the required equations hold when interpreted in $T'(A)$.
%
The proofs of the two laws from the reader theory are analogous to those for the
ordinary reader monad and hence omitted.
%
Next, we consider the tensor equations stating that $\get$ commutes with each
operation in $\calA$. The proof is analogous to the corresponding proof for the
global state monad, and hence omitted.

To prove that the equations in $\calAg$ hold, we first prove a following lemma,
analogous to the one used in the proof for the global state monad transformer.
%
\begin{lemma}
    Let $f : \Gamma \to T'(A)$ be a function. Then for all $r \in R$, there
    exists a function $\fhat_r : \Gamma \to T(R)$ such that, for all terms $t$
    in $F_\Sigma(\Gamma)$, we have
    %
    \[ \sem{t}_{\fhat_r}^{T(A)} = (\sem{t}_f^{T'(A)})(r). \]
\end{lemma}
The proof is similar to that of the analogous lemma for global state, and hence
omitted.

Now we can show that the equations in $\calAg$ hold:
Let $e = (\Gamma, t_1, t_2) \in
E$, and let $f : \Gamma \to T'(A)$ be a function. We need to show

\[ \sem{t_1}_f^{T'(A)} = \sem{t_2}_f^{T'(A)}, \]

where both sides are functions $R \to T(A)$. To this end, let $r \in R$.
Then by the above Lemma, we have that there is $\fhat_r : \Gamma \to T(A)$
with the property described in the Lemma. In particular, we have
%
\begin{align*}
    (\sem{t_1}_f^{T'(A)})(s) 
    &= \sem{t_1}_{\fhat_s}^{T(S \times A)} \\
    &= \sem{t_2}_{\fhat_s}^{T(S \times A)} \\
    &= (\sem{t_2}_f^{T'(A)})(s),
\end{align*}
%
where the second equality holds because $T(A)$ is a model of $\calAg$
and hence the equation holds when interpreted in $T(A)$.

Next, we define $\eta' : A \to T'(A)$ by 
%
$\eta'(x) = \lambda r.\eta(x)$, 
%
where $\eta : A \to T(A)$ is the unit for $T$.

Finally, we show the universal property.
%
Let $B$ be a model of the theory $\calAg \otimes \reader{R}$, and let $f : A \to UB$.
%
We define $\fbar : T'(A) \to B$ as the following composition:
%
\[ \fbar = 
    R \to T(A) \xrightarrow{R \to f^*} R \to B \xrightarrow{B.\sem{\lookup}} B.  \]
%
where $f^*$ is the unique extension of $f$ to a homomorphism of $\calAg$-models.

Then it is easily checked that $\fbar$ is a homomorphism of models of $\calAg
\otimes {\reader{R}}$, and that $\fbar \circ \eta' = f$.

The proof of uniqueness is similar to the corresponding proof for global state,
and is hence omitted.
\end{proof}

Next, we show that the monad transformer lifts to relations:
\begin{lemma}\label{lem:reader-monad-transformer-free-domains}
    Let $\calA$ be an ordinary algebraic theory. Let $\calA^G \in \calA[\mho,
    \theta]$, be a guarded extension of $\calA$. Suppose that the adjunction
    between sets and models of $\calA^G$ lifts to a strong adjunction between
    predomains and $\calA^G$-domains, with the free domains being given by
    $\That$.
    
    Then for the theory $\calA^{G\otimes} := (\calA^{G} \otimes \reader{R})$,
    the adjunction between sets and models of $\calA^{G\otimes}$ lifts to a
    strong adjunction between predomains and $\calA^{G\otimes}$-domains. The
    free domains for $\calA^{G\otimes}$ are given by $R \to \That(A)$.
\end{lemma}
\begin{proof}
    % We define $\calA^{G+}$ to be $\calA^{G} + \mho$. % i.e., exceptions
    Let the free models of $\calA^G$ be given by $T$, and let the free domains
    of $\calA^G$ be given by $\That$. By assumption, for any predomain $A$, we
    have $|\That(A)| = T(|A|)$.
    
    We know by the above that at the level of sets, the free models of
    $\calA^{G+}$ are given by $T'(A) = R \to T(A)$.
    
    We show that $\That'(A) := R \to \That(A)$ gives the free domains of
    $\calA^{G+}$. Here, $R$ is a discrete predomain, so the relations are given
    by equality. In addition, $\to$ denotes the predomain of morphisms of
    predomains between $R$ and $\That(A)$ (this is the same as ordinary
    functions from $R$ to $\That(A)$), and $\times$ denotes the product of
    predomains.

    \begin{itemize}
        \item The strong error ordering and weak bisimilarity on $\That'(A)$ are
        those of the predomain $R \to \That(A)$.

        \item For a predomain relation $c : A \rel A'$, the lifting is given by
        $\That'(c) = R \to \That(c)$.
        %
        More concretely, we have $f \binrel{\That'(c)} g$ iff for all $r \in R$,
        we have $(f\, r) \binrel{\That(c)} (g\, r)$.

        \item To verify the universal property of the lifting of a relation, we
        use the definition of the unique extension of a function $A \to UB$ to a
        homomorphism out of $T'(A)$, as defined in the proof of Lemma
        \ref{lem:reader-monad-transformer-free-models} above. The universal
        property for the lifted relation then follows using the actions of $\to$
        and on squares as well as the universal property of $\That$ on
        relations.

        \item It is easily checked that the operations are monotone and
        bisimilarity-preserving, and that the strength $t : A \times T'(A') \to
        T'(A \times A')$ is a morphism of predomains.
        
        \item Lastly, for any predomain $A$, we have
        %
        \begin{align*}
            |\That'(A)| &= |R \to \That(A)| \\
                        &= R \to |\That(A)| \\
                        &= R \to T(|A|)  \\
                        &= T'(|A|).
        \end{align*}
    \end{itemize}
\end{proof}

Now we are ready to define the action of the reader transformer on graduality
schemata. Let $\calA$ be an ordinary algebraic theory, and let $\calG$ be a
graduality schema for $\calA$.
%
We write $\calAs, \calAsb$ for the guarded extensions of $\calA$ associated with
$\calG$, and we write $T$ for the free-model monad for $\calAs$ on Set, and
$\That$ for its lifting to predomains.
%
We define a graduality schema for $(\calA \otimes \reader{R})$, as follows.
\begin{itemize}
    \item We define the guarded extensions 
    $\calAgtimes := \calAs \otimes \reader{R}$ and 
    $\calAgtimes_b := \calAsb \otimes \reader{R}$.

    \item Then by Lemma \ref{lem:reader-monad-transformer-free-models}, the free models of
    $\calAgtimes$ on Set are given by $T'(A) := R \to T(A)$.

    \item By Lemma \ref{lem:reader-monad-transformer-free-domains}, we have that 
    the free domains for $\calAgtimes$ are given by $R \to \That(A)$.

    \item We define the extensional types $E'(A) := R \to E(A)$, where
    $E$ is the family of extensional types for $\calG$. Next, we define the
    family of functions
    %
    $\col'_A : T'^{gl}(A) \to E'(A)$, i.e., $(R \to T^{gl}(A)) \to (R \to E(A))$, via 
    %
    \[ \col'_A(f) := \col_{A} \circ f, \]
    %
    where $\col$ is the family of functions associated with the schema $\calG$.

    \item Lastly, the relation $\ltbisim_{E'(A)}$ is defined to be $R \to
    \ltbisim_{E(A)}$, i.e., we have $f \binrel{\ltbisim_{E'(A)}} g$ iff
    for all $r \in R$, we have $(f\, r) \binrel{\ltbisim_{E(A)}} (g\, r)$.
    
    Then we observe that $\ltbisim_{E'(A)}$ is a partial order for each $A$. We
    now show that the error ordering and bisimilarity relations on the
    globalized monad imply $\ltbisim_{E'(A)}$ on the output of $\col'_A$.
    %
    More concretely, let $f, g : R \to T^{gl}(A)$, and suppose that 
    $f \ltdyn_{R \to T^{gl}(A)} g$. Then we want to show that
    $\col'_A(f) \ltbisim_{R \to E(A)} \col'_A(g)$. Let $r \in R$.
    It suffices to show
    %
    \[ \col_{A}(f(r)) \ltbisim_{E(A)} \col_{A}(g(r)). \] 
    %
    This follows from the fact that $f(r) \ltdyn_{T^{gl}(A)} g(r)$.

    The same reasoning holds for bisimilarity.
\end{itemize}












